\documentclass[12pt, a4paper]{article}
%%%%% 1. CORE CONFIGURATION
\usepackage{expl3,xparse}
\usepackage[utf8]{inputenc}
\usepackage[T1]{fontenc}
\usepackage{lmodern}
\usepackage{microtype}
\usepackage[table,dvipsnames]{xcolor}
\usepackage{geometry}
\geometry{a4paper,total={170mm,257mm},left=20mm,top=20mm,marginparwidth=20mm}

%%%%% 2. UTILITIES
\usepackage{etoolbox}
\usepackage{settobox}
\usepackage{enumitem}
\usepackage{booktabs}
\usepackage{array}
\usepackage{longtable}
\usepackage{float}
\usepackage{abstract}
\usepackage{fancyhdr}
\usepackage[textwidth=16.3mm,textsize=scriptsize]{todonotes}
    \newcommand*{\newcommentator}[2][red!20]{%
        \expandafter\newcommand\csname #2\endcsname[1]{\todo[color=#1]{{\bf\makefirstuc{#2}:} ##1}}%
        \expandafter\newcommand\csname big#2\endcsname[1]{\todo[inline,color=#1]{{\bf\makefirstuc{#2}:} ##1}}}
    \newcommentator{pablo}
    \newcommentator[blue!20]{arka}


%%%%% 3. GRAPHICS
\usepackage{graphicx}

%%%%% 4. MATH & FONTS
\usepackage{amssymb,amsmath}
\numberwithin{equation}{section}
\usepackage{mathtools}
\usepackage{mleftright}
\mleftright
\usepackage{stackrel}
\usepackage{bm}
\usepackage{bbm}

%%%%% 5. ALGORITHMS
\usepackage[linesnumbered,ruled]{algorithm2e}

%%%%% 6. PROJECT MATH ENVIRONMENTS AND MACROS
\usepackage{preamble/mathenv}
\usepackage{preamble/project-macros}

%%%%% 7. BIBLIOGRAPHY AND LINKS
\usepackage[numbers]{natbib}
\usepackage[pdfencoding=auto]{hyperref}
\hypersetup{
  bookmarksopen = true,
  bookmarksnumbered = true,
  linktoc = all,
  hypertexnames = false,
  hidelinks,
  pdftitle = {AI and Poker: From Reinforcement Learning to Nash Equilibrium},
  pdfauthor = {},
}
\def\figureautorefname{Figure}

%%%%% 8. THEOREM ENVIRONMENTS
\usepackage{preamble/thmenv}
\usepackage[nameinlink]{cleveref}
\crefname{algocf}{algorithm}{algorithms}
\Crefname{algocf}{Algorithm}{Algorithms}

%%%%% 9. ALGORITHM STYLE
\SetAlCapSkip{0.8em}
\SetKwInput{KwIn}{Input}
\SetKwInput{KwOut}{Output}
\SetKwInput{KwInit}{Initialise}
\SetKwInput{KwHyper}{Hyperparameters}
\SetKwInput{KwParam}{Parameters}
\SetKwInput{KwCost}{Cost}
\newcommand{\algcall}[2]{\hyperref[#1]{\texttt{#2}}}
\newcommand\mycommfont[1]{\footnotesize\ttfamily\textcolor{blue!70!black}{#1}}
\SetCommentSty{mycommfont}

%%%%% 10. PAGE STYLE
\setlength{\headheight}{14.5pt}
\pagestyle{fancy}
\fancyhf{}
\fancyhead[L]{\small AI and Poker}
\fancyhead[R]{\small\thepage}
\renewcommand{\headrulewidth}{0.4pt}

\title{%
  \LARGE\bfseries AI and Poker:\\
  From Reinforcement Learning to Nash Equilibrium\\[0.8em]
  \large A Mathematical Synthesis
}
\author{}
\date{May 2026}


\title{Mathematical Notes}
\author{Pablo, Arka}
\date{\today}

\begin{document}
\maketitle

\begin{abstract}
First draft of the research papers to emerge from this project: Rao--Blackwellization, branching and poker case-study.
\end{abstract}

\tableofcontents
\newpage

\section{Introduction}\label{sec:intro}

In the context of end-to-end policy learning, we train a model to learn a parametrized strategy profile $\pi^\theta$.
All policy-gradient methods to train such  $\theta$ are based on a similar principle.
They create a sequence $\theta_0, \ldots \theta_k$ of parameters where typically $\theta_0$ is a random seed, and they gradually build the next iterate by:  
\begin{enumerate}[label =(\roman*)]
  \item \textbf{Rollout:} Sampling games using the current strategy profiles $(\pi_{\theta_i})_{i\le k}$.
  \item \textbf{Loss design:} Secretly, unbeknownst to the opponents, some new player with strategy $\pi^\theta$ will randomly replace the acting player.
  We pick some sort of loss function $\mathcal{L}_{k+1}(\theta)$ which depends on $(\theta_i)_{i\le k}$, and which evaluates how well the new player $\pi^\theta$ is doing.
  \item \textbf{Gradient Descent}: Using minibatches sampled from the online rollout database, we compute some estimator of the loss $\widehat{\mathcal{L}}_{k+1}(\theta)$ and descend along it.
\end{enumerate}

Trajectories in games typically only get assigned to a computable reward at the end of the round (terminal trajectories), which forces rollout to \textit{sample full runouts}.
Typically, each runout is sampled by generating an explicit sequence of full game states (all private and public information is decided).

We introduce here two methods to reduce the variance of the loss estimators $\widehat{\mathcal{L}}_{k+1}(\theta)$:
\begin{enumerate}[label =(\roman*)]
  \item \textbf{Rao--Blackwellization of private state:} games have typically a very wide variance coming from the unavailable private information: the exact same strategy can lead to terrible failure or wonderful success depending on the private state of the opponents.
  We suggest addressing this by integrating away this variance.
  We replace private-state sampling in the loss estimator by conditional expectations over the posterior distribution of private states.

  \item \textbf{Variance-guided Branching:} a significant amount of variance in the loss estimator typically also comes from the variety of public events (actions and chance events) which can significantly modify the strength of the acting player's private state.
  We suggest addressing this focusing the sampling around higher variance spots, by seeing rollouts as a tree instead of a set of sequences.
\end{enumerate}

\section{Game-Theoretical Model}\label{sec:pbg}

To simplify the notation, we focus on \textbf{finite Public-action Bayesian Games (PBG)}: all private information is sampled initially, after which every player action and chance outcome is public and play is sequential.
Player $p$ therefore acts with information state $(x_p,\sigma)$, and these games have perfect recall.
This restricted extensive-form Bayesian class includes many poker, bidding, and auction models, but not dynamic private observations, hidden actions, or simultaneous moves without further extensions.

\begin{definition}[Finite public-action Bayesian game]\label{def:finite-pbg}
For a finite set $E$, let
\begin{equation}\label{eq:probability-simplex}
\Delta(E):=\left\{\lambda:E\to[0,1]\;\middle|\;\sum_{e\in E}\lambda(e)=1\right\}.
\end{equation}

A finite public-action Bayesian game consists of the following data.
\begin{enumerate}[label =(\roman*)]
    \item \textbf{Tree:} A finite prefix-closed set $\mathcal S$ of words, seen as a tree, with empty word $\varnothing\in\mathcal S$.
    Its elements are called public states, or trajectories.
    Write $\sigma\preccurlyeq\tau$ if $\sigma$ is a prefix of $\tau$, and $\sigma\prec\tau$ if $\sigma\preccurlyeq\tau$ and $\sigma\neq\tau$.
    The terminal trajectories are the maximal elements of $\mathcal S$ and are denoted by $\mathcal T$.
    For every nonterminal trajectory $\sigma\in\mathcal S\setminus\mathcal T$, there is a finite event set $\mathcal E_\sigma$ such that the immediate successors of $\sigma$ are precisely the trajectories $\sigma e$, with $e\in\mathcal E_\sigma$.

    \item \textbf{Owner:} A finite set $\mathcal{P}$ of \textit{players} and an owner map $o:\mathcal S\setminus\mathcal T \to \mathcal{P}\sqcup\{\mathsf c\}$.
    If $o(\sigma)\in\mathcal P$ (resp. $\{\mathsf c\}$) is a player node (resp. chance node), we typically denote the elements of $\mathcal{E}_\sigma$ by $a$ (resp. $c$).
    For $\sigma \in \mathcal S$ and $p \in \mathcal P$, we write
    $|\sigma|_p:=\#\{\rho\preccurlyeq\sigma\mid o(\rho)=p\}$ for the number of player nodes in $\sigma$ owned by $p$, and $|\sigma|_{\textnormal{player}}:=\sum_{p\in\mathcal P}|\sigma|_p$ for the total number of player nodes in $\sigma$.
    Thus, when $o(\sigma)=p$, the state $\sigma$ is counted in $|\sigma|_p$: it is the state immediately before $p$ chooses an outgoing action $a\in\mathcal E_\sigma$; that action is the edge from $\sigma$ to the next state $\sigma a$.
    For $\rho\preccurlyeq\sigma$, we write $|\sigma\setminus\rho|_p:=|\sigma|_p-|\rho|_p$ and $|\sigma\setminus\rho|_{\textnormal{player}}:=|\sigma|_{\textnormal{player}}-|\rho|_{\textnormal{player}}$.
    For $\tau \in \mathcal{T}$, write $\tau^{\mathsf{pre-c}}$ for the largest prefix ending in no more than one node in $\mathcal{T} \sqcup o^{-1}(\mathsf{c})$.
    In other words, we have $\tau^{\mathsf{pre-c}} = \sigma a$ for $\sigma a \preccurlyeq \tau$ the last player node $o(\sigma) \in \mathcal{P}$ and decision $a$: if $\sigma a \prec \tau$, we must have $o(\rho) = \mathsf{c}$ for every $\sigma a \preccurlyeq \rho \prec \tau$.
    Informally, a pre-terminal state is when all players have acted, and before any final chance events are revealed which may impact rewards (\textit{e.g.} when players are all-in in poker before the board is complete, the pre-terminal state is right after the last ``\textit{call}'' and before seeing any more board cards).

    \item \textbf{Initial Private State Profile:} For every player $p\in \mathcal P$, a finite type set $\mathcal X_p$, and a distribution $\mu$ on $\mathcal X:=\prod_{p\in\mathcal P}\mathcal X_p$.
    Elements of $\mathcal X$ are typically written $x= x_{\mathsf c} = (x_p)_{p \in \mathcal P}$, with $x_p\in\mathcal X_p$.

    \item \textbf{Public-Chance Rules:} For every chance node $\sigma$ and every initial profile $x \in \mathcal{X}$, a public-chance rule $\chi(\cdot \mid x, \sigma)\in \Delta(\mathcal E_\sigma)$.

    \item \textbf{Reward:} a function $\mathbf{R} = (R_p)_{p \in \mathcal{P}}: \mathcal{X} \times \mathcal{T}\to\R^{\mathcal{P}}$.
\end{enumerate}
\end{definition}

\begin{definition}[Strategy and Random Trajectories]\label{def:strategy-random-trajectories}
A strategy profile is the data of $\pi(\cdot \mid x_p, \sigma) \in \Delta(\mathcal E_\sigma)$ for $\sigma \in o^{-1}(\mathcal{P})$ and $x_{o(\sigma)} \in \mathcal{X}_{o(\sigma)}$.
For $p \in \mathcal{P}$, we write $\pi_p$ for the strategy of player $p$, \textit{i.e.} the subcollection of these distributions for $p = o(\sigma)$, and $\pi_{-p}$ for the strategies of all other players.
A strategy profile induces a joint probability measure $\mathbb{P}^\pi$ on the set $\mathcal{X} \times \mathcal{T}$.

For a trajectory $\sigma\in\mathcal S$, a player $p$ and a private initial profile $x_p \in \mathcal{X}_p$, define the \textbf{reach} as the probability that player $p$ acted according to $\sigma$ conditional on initial profile $x_p$ and every other event following $\sigma$:
\begin{equation}\label{eq:player-reach}
  r_p^\pi(\sigma\mid x_p)
  := \mathbb{P}^\pi(p \textnormal{ acts as }\sigma \mid x_p, \textnormal{other events follow }\sigma)
  = \prod_{\substack{\rho a \preccurlyeq \sigma \\ o(\rho) = p}} \pi(a \mid x_p, \rho),
\end{equation}
where the product ranges over the actions of player $p$ in the trajectory $\sigma$.
Similarly, we define the \textbf{draw odds} of a trajectory $\sigma$ as the probability that the public chance events in $\sigma$ occurred conditional on the initial profile $x$:
\begin{equation}[rCl]
  r_{\mathsf c}(\sigma\mid x)
  &:=& \mathbb{P}^\pi(\textnormal{chance happens as in }\sigma \mid x, \textnormal{players act as } \sigma)
  \nonumber\\
  &=& \prod_{\substack{\rho c \preccurlyeq \sigma \\ o(\rho) = \mathsf c}} \chi(c \mid x, \rho),
  \label{eq:chance-reach}
\end{equation}
The expected payoff of player $p$ is
\begin{equation}\label{eq:expected-payoff}
  J_p(\pi) :=\mathbb{E}^\pi(R_p) = 
  \mathbb{E}_{X \sim \mu} 
  \left(
      \sum_{\tau\in\mathcal T} R_p(X, \tau) \prod_{q \in \mathcal{P} \cup \{\mathsf c\}} r_q(\tau \mid X_q)
  \right)
\end{equation}
For $\sigma \in o^{-1}(\mathcal P)$, $p = o(\sigma)$ and $x_p \in \mathcal{X}_p$, we write $V(x_p, \sigma) = \mathbb{E}^\pi(R_{p}(x, \tau) \mid x_p,\sigma)$.
For a further choice of $a \in \mathcal{E}_\sigma$, we write $Q(x_p, \sigma a) = \mathbb{E}^\pi(R_p(x, \tau) \mid x_p, \sigma a)$ and $A(x_p, \sigma a) = Q(x_p, \sigma a) - V(x_p, \sigma)$.
Here and throughout, a conditional expectation conditioned on an event of probability zero is defined to be zero.
\end{definition}

\begin{remark}[Evolving private states]
One could allow the private state to evolve by assigning
$\mathcal X_\sigma=\prod_p\mathcal X_{\sigma,p}$ to every public state.  At a
player node $o(\sigma)=p$, for each $x_p\in\mathcal X_{\sigma,p}$ let
\[
  \mathcal D_\sigma(x_p)\subseteq
  \coprod_{a\in\mathcal E_\sigma}
  \{a\}\times\mathcal X_{\sigma a,p}
\]
be a nonempty set of legal decisions.  The strategy would be a transition
kernel
\[
  \pi(\cdot\mid x_p,\sigma)\in\Delta(\mathcal D_\sigma(x_p)),
\]
so that the player jointly chooses the public action and their next private
state.  At a chance node, we would similarly take
\[
  \chi(\cdot\mid x,\sigma)\in
  \Delta\!\left(\coprod_{c\in\mathcal E_\sigma}
  \{c\}\times\mathcal X_{\sigma c}\right).
\]
Writing $\mathbf r^\pi(\sigma)(x)$ for unnormalised reach mass, and carrying
the nonacting coordinates unchanged, reach is defined inductively by
\[
\mathbf r^\pi(\varnothing)(x)=\mu(x),\qquad
\mathbf r^\pi(\sigma e)(x')=
\begin{cases}[l"l]
\displaystyle\sum_{x:\,x_{-p}=x'_{-p}}
\mathbf r^\pi(\sigma)(x)\pi(e,x'_p\mid x_p,\sigma)
& o(\sigma)=p,\\
\displaystyle\sum_{x\in\mathcal X_\sigma}
\mathbf r^\pi(\sigma)(x)\chi(e,x'\mid x,\sigma)
& o(\sigma)=\mathsf c.
\end{cases}
\]
Thus $\mathbb P^\pi(\sigma\preccurlyeq\tau)=\sum_x\mathbf r^\pi(\sigma)(x)$,
and the posterior at $\sigma$ is its normalisation.  These kernels give a much
larger class of games to which the same techniques apply.  We keep the
formalism below closer to poker, where private states are sampled initially,
to simplify the notation.
\end{remark}

\begin{notation}[Stationary Distributions in Tree-Shaped Markov Chains]\label{not:stationary-distribution}
  Whenever we have a tree-shaped Markov chain with a collection of terminal states, we call \textit{stationary distribution} the unique probability distribution that is stable under:
  \begin{enumerate}[label = (\roman*)]
    \item at any non-terminal state, randomly moving down the tree according to the transition probabilities, and
    \item at any terminal state, jumping back to the root of the tree.
  \end{enumerate}
  From a game persepective, we see games as endlessly restarting whenever someone won.
\end{notation}

\begin{notation}[Ideal Random Variables]\label{not:random-variables}
We first sample an ideal player-decision node $\Sigma^\pi$ according to
\begin{equation}\label{eq:stationary-node-law}
  \mathbb{P}(\Sigma^\pi = \sigma)
  = \one_{o(\sigma)\in\mathcal P}
  \frac{\mathbb{P}^\pi(\sigma \preccurlyeq \tau)}
       {\mathbb{E}^\pi(|\tau|_{\textnormal{player}})}.
\end{equation}
Conditionally on $\Sigma^\pi=\sigma$, we let $X^\pi$ follow the posterior law
\begin{equation}\label{eq:ideal-private-profile-law}
  \mathbb{P}(X^\pi=x\mid\Sigma^\pi=\sigma)
  =
  \frac{\mu(x)\mathbb{P}^\pi(\sigma\preccurlyeq\tau\mid X=x)}
       {\mathbb{P}^\pi(\sigma\preccurlyeq\tau)},
\end{equation}
we set $P^\pi=o(\Sigma^\pi)$, and we let
$A^\pi\sim\pi(\cdot\mid X^\pi_{P^\pi},\Sigma^\pi)$.
Before observing $\Sigma^\pi$, the joint and marginal laws of $X^\pi$ and $P^\pi$ are
\begin{equation}\label{eq:ideal-profile-player-law}
  \mathbb{P}(X^\pi=x,P^\pi=p)
  = \frac{\mu(x)\mathbb{E}^\pi(|\tau|_p\mid X=x)}
           {\mathbb{E}^\pi(|\tau|_{\textnormal{player}})},
\end{equation}
and hence
\begin{equation}\label{eq:ideal-profile-and-player-marginals}
  \mathbb{P}(X^\pi=x)
  = \frac{\mu(x)\mathbb{E}^\pi(|\tau|_{\textnormal{player}}\mid X=x)}
           {\mathbb{E}^\pi(|\tau|_{\textnormal{player}})},
  \qquad
  \mathbb{P}(P^\pi=p)
  = \frac{\mathbb{E}^\pi(|\tau|_p)}
           {\mathbb{E}^\pi(|\tau|_{\textnormal{player}})}.
\end{equation}
When $\pi$ is clear from the context, we drop the superscript and write $X$, $P$, $\Sigma$ and $A$.
\end{notation}

Recall below how the key theorem of policy gradient theory adapts to this setting.

\begin{theorem}[Policy Gradient Theorem for Games]\label{thm:policy-gradient-games}
  For $p \in \mathcal{P}$, using the ideal random variables of \autoref{not:random-variables}, we have
  \begin{equation}\label{eq:policy-gradient-games}
    \nabla_{\pi_p} J_p(\pi)
    = \mathbb{E}^\pi(|\tau|_{p}) \nabla_{\pi_p} \mathbb{E}^{\sg[\pi]}\left(\log(\pi(A \mid X_p, \Sigma)) \sg[Q(X_p, \Sigma A)] \mid P = p\right)
  \end{equation}
\end{theorem}
Thus the aggregate profile gradient targeted below is the sum of the unilateral gradients in \autoref{eq:policy-gradient-games}; equivalently, it weights each player's conditional gradient by $\mathbb P(P=p)=\mathbb E^\pi(|\tau|_p)/\mathbb E^\pi(|\tau|_{\textnormal{player}})$ under the ideal law.

\section{Self-Play Policy-Gradient Methods in Synchronous Rollouts}\label{sec:synchronous-pg}

\subsection{Theoretical Model}\label{subsec:pg-theoretical-model}

Assuming a \textit{synchronous rollout}, all main policy-gradient methods to solve games can be modelled by the following choices:
\begin{enumerate}
  \item A \textbf{model} is the choice of a parametrized strategy profile $\theta \mapsto \pi^\theta$.
  Optionally, it also included the choice of a parametrized value functions $\theta \mapsto V^\theta$ meant to approximate $\mathbb{E}^{\pi^\theta}(R_p(x,\tau) \mid x_p, \sigma \preccurlyeq \tau)$ for all $\sigma \in \mathcal{S}$, $p = o(\sigma)$ and $x_p \in \mathcal{X}_p$.

  \item Some \textbf{value} estimators 
  \begin{equation}[rCl]
    \widehat{A}(x_p, \sigma a) &\simeq& \mathbb{E}^\pi(R_{p}(x,\tau) \mid x_p, \sigma a) - \mathbb{E}^\pi(R_{p}(x,\tau) \mid x_p, \sigma)
    \\
    \widehat{Q}(x_p, \sigma a) &\simeq& \mathbb{E}^\pi(R_{p}(x,\tau) \mid x_p, \sigma a)
    \\
    \widehat{V}(x_p, \sigma) &\simeq& \mathbb{E}^\pi(R_{p}(x,\tau) \mid x_p, \sigma)
  \end{equation}
  for $\sigma a \in \mathcal{S}$, $p = o(\sigma)$ and $x_p \in \mathcal{X}_p$. 
  Optionally, the reward $R_p$ above can be replaced by a modified reward depending on $\pi$ and some reference profile $\rho$ (see R-NaD).
  
  \item A \textbf{loss function} $\mathcal{L}$ from which we want to estimate
  \begin{equation}\label{eq:pg-loss-expectation}
    l(\theta) = \mathbb{E}\left(
      \mathcal{L}\left(
        (\pi^\theta, V^\theta,\pi, V)(X_P, \Sigma),
        A,
        \widehat{A},
        \widehat{Q},
        \widehat{V}
      \right)
    \right)
  \end{equation}
  for $\Sigma, A, P, X$ as in \autoref{not:random-variables} and $V$ all based on a frozen policy $\pi$, called the \textit{rollout profile} (typically an older version of $\pi^\theta$), and for $\widehat{A} = \widehat{A}(X_P, \Sigma A)$, $\widehat{Q} = \widehat{Q}(X_P, \Sigma A)$ and $\widehat{V} = \widehat{V}(X_P, \Sigma)$.
  Optionally, the loss function may also depend on $\alpha(\cdot \mid X_P, \Sigma)$ for reference profile $\alpha$ (see NashPG).
\end{enumerate}

The \textbf{training} is then an inductive process defining a sequence of weights $(\theta_k)_{k \geq 0}$.
The first choice $\theta_0$ is a random seed sampled from some initial distribution.
Assuming we constructed $\theta_0, \ldots, \theta_k$, we define the next iterate $\theta_{k+1}$ by: estimating $l(\theta)$ using the rollout profile $\theta_k$, then performing some number of gradient descent steps on $\theta$ to minimize the estimated loss.

\paragraph{Assumptions.}
In the model above, we assume a few things that we justify here:
\begin{enumerate}
  \item \textbf{No $Q$-head:} In many games of interests, the action space is too large to have a reasonnable $Q$-head as output of the model.
  \item \textbf{No other-player value head:} With imperfect-information games, the state of the other player is not known.
  To output a value function for each player, one would need to $\sum_p |\mathcal{X}_p|$ outputs, which is typically too large to be practical.
  It's also not a quantity that we necessarily want the model to spend effort to learn: the expected returns over other players \textbf{without} conditioning on their private state is a much more relevant quantity to learn (see \autoref{eq:V-vector-head}).
  \item \textbf{Synchronous rollout:} This is purely a theoretical assumption to simplify the exposition.
\end{enumerate}
These assumptions shape the estimators and losses that are relevant for our context.

\paragraph{Primitives.}
For $0 \leq i \leq k$, we let $\rho^\theta_i = \frac{\pi^\theta(A \mid X_{P},\Sigma)}{\pi^{\theta_i}(A \mid X_{P},\Sigma)}$.
We define \textbf{centered logits} as
\begin{equation}\label{eq:centered-logits}
  z^\theta(a \mid X_{P}, \Sigma) := \log\pi^\theta(a \mid X_{P},\Sigma)
  -
  \frac{1}{|\mathcal E_\Sigma|}
  \sum_{b\in\mathcal E_\Sigma}
  \log\pi^\theta(b\mid X_{P},\Sigma)
\end{equation}
Write $\widehat{A} = \widehat{Q}(X_{P},\Sigma A) - V^{\theta_k}(X_{P},\Sigma)$.
If we are able to compute one, for $b \in \mathcal{E}_\Sigma$, write $\widehat{Q}(b) = \widehat{Q}(X_{P}, \Sigma b)$ an estimator of an all-action target.
We also define a clipped value function $V_{\tt clip}^\theta = V^{\theta_k} + \clip(V^\theta - V^{\theta_k}, -\varepsilon_V, \varepsilon_V)$.
We then have the following classic losses (the notation $\mathcal{L}[\widehat{E}]$ means ``the loss $\mathcal{L}$ depends on a choice of estimator $\widehat{E}$''):
\begin{equation}[rCl]
  \mathcal L_{\tt PG}[\widehat{A}]
  & := &
  -\rho_k^\theta \widehat{A},
  \nonumber\\
  \mathcal L_{\tt clip}[\widehat{A}]
  & := &
  -\min\left\{
    \rho_k^\theta \widehat{A},
    \clip\left(
      \rho_k^\theta, 1-\varepsilon,1+\varepsilon
    \right)
    \widehat{A}
  \right\}
  \nonumber\\
  \mathcal L_{\tt NeuRD}[\widehat{Q}(-)]
  & := &
  - \sum_{b \in \mathcal{E}_\Sigma}
  z^\theta( b \mid X_{P},\Sigma) \cdot
  \clip\left(
    \widehat{Q}(b),
    -c_{\tt NeuRD},c_{\tt NeuRD}
  \right)
  \nonumber\\
  \mathcal L_{\tt NeuRD}^{\tt sample}[\widehat{Q}]
  & := &
  - \frac{z^\theta( A \mid X_{P},\Sigma)}{\pi^{\theta_k}(A \mid X_{P},\Sigma)} \cdot
  \clip\left(
    \widehat{Q},
    -c_{\tt NeuRD},c_{\tt NeuRD}
  \right)
  \nonumber\\
  \mathcal L_{\tt V}[\widehat{V}]
  & := &
  c_V\left(V^\theta(X_{P},\Sigma) - \widehat{V}\right)^2,
  \nonumber\\
  \mathcal L_{\tt V-clip}[\widehat{V}]
  & := &
  c_V \max\left\{\left(V^\theta(X_{P},\Sigma) - \widehat{V}\right)^2, \left(V_{\tt clip}^\theta(X_{P},\Sigma) - \widehat{V}\right)^2\right\},
  \nonumber\\
  \mathcal L_{\tt H}
  & := &
  -c_H\Ent\!\left(\pi^\theta(\cdot\mid X_{P},\Sigma)\right),
  \nonumber\\
  \mathcal L_{\tt KL}
  & := &
  c_{\tt KL} \KLop\left(\pi^\theta(\cdot\mid X_{P},\Sigma) \middle\| \pi^{\theta_{\text{ref}(k)}}(\cdot\mid X_{P},\Sigma)\right)
  \nonumber\\
  \mathcal L_{\tt KL}^{\tt sample}
  & := &
  c_{\tt KL}\left(\rho^\theta_{k} \log(\rho^\theta_{k}) - \rho^\theta_{k} + 1\right)
  \label{eq:primitive-losses}
\end{equation}
The displayed NeuRD losses clip raw $Q$-values and are therefore not invariant to additive reward shifts; they should be understood as deliberate non-invariant surrogates.

\paragraph{Relation to \autoref{thm:policy-gradient-games}.}
The goal of this loss function is that its gradient resembles the gradient of the expected payoff $J_p(\pi^\theta_p, \sg[\pi^\theta_{-p}])$.

\paragraph{Practicality.}
In practice, it's essentially impossible to sample $\Sigma$ in a non-biased way (one needs to first know the average length of the game).
It's also impossible to create the estimators $\widehat{Q}$ and $\widehat{V}$ without seeing the sampled $\Sigma$ as part of a full trajectory.
These two constraints pushes us towards the approximations given in the following subsections.

\subsection{Usual Implementation}\label{subsec:pg-usual-implementation}

We explain here the usual implementation of a training algorithm.

\begin{algorithm}[H]
\caption{$\mathfrak{T}\gets$~\texttt{SampleTerminal}$(\pi,N)$}
\label{alg:sample-terminal}
\DontPrintSemicolon
\tcc{Terminal trajectory sampling used by ordinary sampled-private-state training.}
\KwParam{$\pi$, rollout profile; $N$, number of independent terminal rollouts}
\KwOut{$\mathfrak{T}=((X_i,\Tau_i))_{i=1}^N$}
\For{$i=1,\ldots,N$}{
  Sample $X_i\sim\mu$\;
  Set $\sigma\gets\varnothing$\;
  \While{$\sigma\notin\mathcal T$}{
    $p \gets o(\sigma)$\;
    Sample
    \[
      E \sim \begin{cases}[l"l]
      \pi(\cdot\mid X_{i,p}, \sigma) & \textnormal{if } o(\sigma) \in \mathcal{P}\\
      \chi(\cdot\mid X_i, \sigma) & \textnormal{otherwise}
    \end{cases}
    \]\;
    Set $\sigma \gets \sigma E$\;
  }
  Set $\Tau_i \gets \sigma$\;
}
\Return{$\mathfrak{T}$}
\end{algorithm}

\begin{definition}[Empirical terminal-rollout decision law]\label{def:terminal-rollout-decision-law}
Fix a terminal rollout batch $\mathfrak{T}=((X_i,\Tau_i))_{i=1}^N$ produced by \cref{alg:sample-terminal}.
Define an index $I\in\{1,\ldots,N\}$ by
\begin{equation}\label{eq:terminal-rollout-index-law}
  \mathbb{P}(I=i\mid\mathfrak{T})
  =
  \frac{|\Tau_i|_{\textnormal{player}}}
       {\sum_{j=1}^N|\Tau_j|_{\textnormal{player}}}.
\end{equation}
Then let $\Sigma A \preccurlyeq \Tau_I$ be a uniform choice of a player decision and $P = o(\Sigma)$.
\end{definition}

Do note that all convergence statements are restricted to the fixed finite game, fixed rollout profile, and the sampling schemes stated in their lemmas.

\begin{lemma}[Convergence of the empirical terminal-rollout decision law]\label{lem:terminal-rollout-decision-law-convergence}
Let $\widehat{\mathbb P}_N$ be the conditional law of $(X_I,P,\Sigma,A)$ given $\mathfrak T$, and let $\mathbb P_\star$ be the ideal law of \autoref{not:random-variables}. Then
\[
  \|\widehat{\mathbb P}_N-\mathbb P_\star\|_{\rm TV}\longrightarrow 0
  \quad\textnormal{almost surely},
  \qquad
  \sqrt N(\widehat{\mathbb P}_N-\mathbb P_\star)\Longrightarrow\mathcal N(0,\Gamma)
\]
for a covariance matrix $\Gamma$. Moreover, there are game-dependent constants $c,C>0$ such that, for every $0<\varepsilon<1$,
\[
  \mathbb P\!\left(\|\widehat{\mathbb P}_N-\mathbb P_\star\|_{\rm TV}>\varepsilon\right)
  \leq C e^{-cN\varepsilon^2}.
\]
\end{lemma}

\paragraph{Primitives.}
For $\lambda \in [0,1]$, $1 \leq i \leq N$, $\sigma a \preccurlyeq \Tau_i$ and $p = o(\sigma)$, we write
\begin{equation}\label{eq:sequential-lambda-return}
  \widehat{Q}^\lambda(X_{i, p}, \sigma a)
  :=
  \lambda^{|\Tau_i \setminus \sigma|_p} R_{p}(X_{i}, \Tau_i)
  + (1 - \lambda) \sum_{\substack{\sigma \prec \rho \preccurlyeq \Tau_i \\ o(\rho) = p}}
  \lambda^{|\rho \setminus \sigma|_p - 1} V^{\theta_k}(X_{i, p},\rho)
\end{equation}
and $\widehat{A}^\lambda = \widehat{Q}^\lambda(X_{I,P},\Sigma A) - V^{\theta_k}(X_{I, P},\Sigma)$ and $\widehat{V}^\lambda = \widehat{Q}^\lambda(X_{I,P},\Sigma A)$.
Note that $\lambda = 1$ recovers the Monte Carlo return estimator (necessary in the absence of a value head).
\begin{lemma}[Player-clock GAE]\label{lem:player-clock-gae}
The return of \autoref{eq:sequential-lambda-return} is the usual GAE lambda return for terminal-only rewards, $\gamma=1$, and the successive decisions of player $p$ as time steps.
\end{lemma}
For $|\mathcal{P}| = 2$, $\eta > 0$ and $(x,\tau) \in \mathcal{X} \times \mathcal{T}$, we define the \textbf{\texttt{R-NaD} modified return} as
\begin{equation}\label{eq:rnad-modified-return}
  R^{\tt NaD}_{p}(x, \tau) :=
  R_p(x, \tau) +\eta \sum_{\substack{\sigma a\preccurlyeq\Tau \\o(\sigma) \in \mathcal{P}}} (1-2\one_{o(\sigma)=p})
  \log\left(
    \frac{\pi^{\theta_k}(a\mid x_{o(\sigma)},\sigma)}
    {\pi^{\theta_{\operatorname{ref}(k)}}(a\mid x_{o(\sigma)},\sigma)}
  \right)
\end{equation}
Using this instead of $R$ in the formula above, we get $\widehat{Q}_{\tt NaD}^{\eta, \lambda}$, $\widehat{A}_{\tt NaD}^{\eta, \lambda}$ and $\widehat{V}_{\tt NaD}^{\eta, \lambda}$.
For $b \in \mathcal{E}_\Sigma$, we define the \textbf{\texttt{R-NaD} all-action return} 
\begin{equation}[rCl]
  \widehat{Q}_{\tt NaD}^{\eta, \lambda}(b)
  =
  V^{\theta_k}(X_{I,P}, \Sigma)
  &-&
  \eta\log
  \left(
  \frac{
    \pi^{\theta_k}(b\mid X_{I,P},\Sigma)
  }{
    \pi^{\theta_{\operatorname{ref}(k)}}(b\mid X_{I,P},\Sigma)
  }
  \right)
  \nonumber\\
  &+&
  \frac{\mathbf{1}_{\{b=A\}}}{\pi^{\theta_k}(b\mid X_{I, P},\Sigma)}
  \left(
    \widehat{A}_{\tt NaD}^{\eta, \lambda}
  \ + \eta\log
  \left(
  \frac{
    \pi^{\theta_k}(b\mid X_{I,P},\Sigma)
  }{
    \pi^{\theta_{\operatorname{ref}(k)}}(b\mid X_{I,P},\Sigma)
  }
  \right)
  \right).
\end{equation}

\paragraph{Choices.} Here are typical choices of losses used in training.
\begin{center}
\setlength{\tabcolsep}{3pt}
\scriptsize
\begin{tabular}{@{}
  >{\raggedright\arraybackslash}p{0.17\linewidth}
  >{\raggedright\arraybackslash}p{0.17\linewidth}
  >{\raggedright\arraybackslash}p{0.17\linewidth}
  >{\raggedright\arraybackslash}p{0.43\linewidth}
@{}}
\toprule
\textbf{method}
& \textbf{\text{ref}(k)}
& \textbf{$V^\theta$}
& \textbf{minimized loss}
\\
\midrule
\tt REINFORCE
& $\varnothing$
& $\varnothing$ or fixed baseline
& $\mathcal L_{\tt PG}[\widehat A^1]$
\\
\tt Actor--critic
& $\varnothing$
& $V^\theta$
& $\mathcal L_{\tt PG}[\widehat A^\lambda]+\mathcal L_{\tt V}[\widehat V^\lambda]+\mathcal L_{\tt H}$
\\
\tt PPO
& $\varnothing$
& $V^\theta$
& $\mathcal L_{\tt clip}[\widehat A^\lambda]+\mathcal L_{\tt V}[\widehat V^\lambda]+\mathcal L_{\tt H}$
\\
{\tt MMD}-style {\tt PPO}
& $\varnothing$
& $V^\theta$
& $\mathcal L_{\tt clip}[\widehat A^\lambda]+\mathcal L_{\tt V}[\widehat V^\lambda]+\mathcal L_{\tt H}+\mathcal L_{\tt KL}^{\tt sample}$
\\
\tt NashPG-PPO
& $K\lfloor k/K\rfloor$
& $V^\theta$
& $\mathcal L_{\tt clip}[\widehat A^\lambda]+\mathcal L_{\tt V}[\widehat V^\lambda]+\mathcal L_{\tt KL}$
\\
\tt RNaD-PPO
& $\operatorname{ref}(k)$
& $V^\theta$
& $\mathcal L_{\tt clip}[\widehat A^{{\tt NaD},\lambda}]+\mathcal L_{\tt V}[\widehat V^{{\tt NaD},\lambda}]$
\\
\tt R-NaD / DeepNash
& $\operatorname{ref}(k)$
& $V^\theta$
& $\mathcal L_{\tt NeuRD}[\widehat{Q}_{\tt NaD}^{\eta, \lambda}(-)]+\mathcal L_{\tt V}[\widehat V^{{\tt NaD},\lambda}]$
\\
\bottomrule
\end{tabular}
\end{center}
The method names in this table denote synchronous loss surrogates, not exact reproductions or convergence claims.
In \texttt{R-NaD}, it is not clear what function $\operatorname{ref}(k)$ to choose, as the original implementation is \texttt{IMPALA}-style rather than synchronous.

\paragraph{Implementation.} Below is a pseudocode of the usual implementation.

\begin{algorithm}[H]
\caption{$(\theta_k)_{k\ge0}\gets$~\texttt{SequentialPG}$(\theta_0)$}
\label{alg:sequential-pg}
\DontPrintSemicolon
\tcc{Sequential rollout policy-gradient template corresponding to the procedure above.}
\KwParam{$\theta_0$, initial weights; $\theta$, working copy of the next iterate}
\KwHyper{$N$, rollout count; $B$, minibatch size; chosen number of gradient-descent steps per iterate}
\KwOut{$(\theta_k)_{k\ge0}$}
\For{$k=0,1,2,\dots$}{
  Sample $(X_i,\Tau_i)_{i=1}^N \stackrel{\mathsf{iid}}{\sim} (X, \Tau^{\pi^{\theta_k}})$ using $\algcall{alg:sample-terminal}{SampleTerminal}(\pi^{\theta_k},N)$\;
  \For{$i=1,\ldots,N$}{
    Compute $R(X_i, \Tau_i)$ (optionally, using \autoref{eq:rnad-modified-return})\;
    By backward propagation, for $\sigma a \preccurlyeq \Tau_i$ and $p = o(\sigma)$, compute $\widehat{Q}(X_{i, p}, \sigma a)$ using \autoref{eq:sequential-lambda-return}\;
    Compute $\widehat{V}(X_{i, p}, \sigma)$\;
  }
  $\theta\gets\theta_k$\;
  \For{each gradient-descent step}{
    Sample $(I_b, \Sigma_b, A_b, P_b)_{b=1}^B \stackrel{\mathsf{iid}}{\sim} (I, \Sigma, A, P)$ as in \autoref{def:terminal-rollout-decision-law}\;
    Deduce $\widehat{A}_b$, $\widehat{Q}_b$, and $\widehat{V}_b$ at $(X_{I_b,P_b},\Sigma_b,A_b)$\;
    Estimate
    \begin{equation*}[rCl]
      \widehat{l}(\theta)
      &=& \frac{1}{B}\sum_{b=1}^B \mathcal{L}\left(
        (\pi^\theta, V^\theta,\pi^{\theta_k}, V^{\theta_k})(X_{I_b, P_b}, \Sigma_b),
        A_b,
        \widehat{A}_b,
        \widehat{Q}_b,
        \widehat{V}_b
        \right)
    \end{equation*}\;
    Descend $\theta$ along $\nabla_\theta \widehat{l}(\theta)$\;
  }
  store $\theta_{k+1}=\theta$\;
}
\Return $(\theta_k)_{k\ge0}$
\end{algorithm}

Optionally, for some $\operatorname{ref}(k) < k$ reference index, during minibatches, we can evaluate $\pi^{\theta_{\operatorname{ref}(k)}}(\cdot \mid X_{I,P},\Sigma)$ at the sampled $(I, P, \Sigma)$ to modify the loss function (see NashPG).

\subsection{Rao--Blackwellization of the Private State}\label{subsec:rao-blackwell-private-state}

\paragraph{Motivation.}
Game strategies often involve disguising one's private information.
In our leading example, poker, good players must ``bluff'' (pretend to have good cards when they don't), ``slowplay'' (pretend to have bad cards when they don't) and more generally, they must ``protect their action ranges'', meaning that $\mathbb{P}(\text{private state of } p\mid \text{public information})$ should not be too revealing.
When learning poker, players are taught to ``balance their ranges'' and think ``what would I do if I had a different hand?''.
Motivated by this, it's a natural choice to design models so that it doesn't see his private state $x_p$ but rather decides what to do with any potential private state $x_p$.
In other words, we do a currying transformation and see the model as
\begin{equation}\label{eq:rao-blackwell-curried-policy}
  \theta \longmapsto \left(x_p \in \mathcal{X}_p \mapsto \pi^\theta(\cdot \mid x_p, \sigma) \in \Delta(\mathcal{E}_\sigma)\right)
\end{equation}
which means that a single model run now outputs a collection of distributions.
The goal of Rao--Blackwellization is to exploit this architecture in the training algorithm to reduce the variance of the loss estimator induced by the private state $X$.

\begin{algorithm}[H]
\caption{$\mathfrak{T}^{\mathsf{pre-c}}\gets$~\texttt{SamplePreTerminal}$(\pi,N)$}
\label{alg:sample-preterminal}
\DontPrintSemicolon
\tcc{Public pre-terminal sampling for Rao--Blackwellized training.}
\KwParam{$\pi$, rollout profile; $N$, number of independent public rollouts}
\KwOut{$\mathfrak{T}^{\mathsf{pre-c}}=((\Tau_i^{\mathsf{pre-c}},\mathfrak r_i))_{i=1}^N$}
\For{$i=1,\ldots,N$}{
  Sample $X_i\sim\mu$\;
  Set $\sigma\gets\varnothing$ and $\mathfrak{r}_i\gets (1)_{q\in\mathcal P\cup\{\mathsf c\}}$\;
  \While{$\sigma\notin\mathcal{T}^{\mathsf{pre-c}}$}{
    Set $q\gets o(\sigma)$\;
    Sample
    \[
      E \sim \begin{cases}[l"l]
      \pi(\cdot\mid X_{i,q}, \sigma) & \textnormal{if } q \in \mathcal{P}\\
      \chi(\cdot\mid X_i, \sigma) & \textnormal{otherwise}
      \end{cases}
    \]\;
    Set
    \[
      \mathfrak{r}_{i,q} \gets \mathfrak{r}_{i,q} \odot \begin{cases}[l"l]
      \pi(E\mid -, \sigma) & \textnormal{if } q \in \mathcal{P}\\
      \chi(E\mid -, \sigma) & \textnormal{otherwise}
      \end{cases}
    \]\;
    Normalize $\mathfrak{r}_{i,q}$ and set $\sigma\gets\sigma E$
  }
  Set $\Tau_i^{\mathsf{pre-c}} \gets \sigma$\;
}
\Return{$\mathfrak{T}^{\mathsf{pre-c}}$}
\end{algorithm}

\begin{definition}[Empirical pre-terminal Rao--Blackwell decision law]\label{def:preterminal-rao-blackwell-decision-law}
Fix a pre-terminal rollout batch $\mathfrak{T}^{\mathsf{pre-c}}=((\Tau_i^{\mathsf{pre-c}},\mathfrak r_i))_{i=1}^N$ produced by \cref{alg:sample-preterminal}.
Define $I\in\{1,\ldots,N\}$ by
\begin{equation}\label{eq:preterminal-rollout-index-law}
  \mathbb{P}(I=i\mid\mathfrak{T}^{\mathsf{pre-c}})
  =
  \frac{|\Tau_i^{\mathsf{pre-c}}|_{\textnormal{player}}}
       {\sum_{j=1}^N|\Tau_j^{\mathsf{pre-c}}|_{\textnormal{player}}}.
\end{equation}
Then let $\Sigma A \preccurlyeq \Tau_I^{\mathsf{pre-c}}$ be a uniform choice of a player decision and $P = o(\Sigma)$.
\end{definition}

\begin{lemma}[Convergence of the empirical pre-terminal decision law]\label{lem:preterminal-decision-law-convergence}
Let $\widehat{\mathbb P}^{\mathsf{pre-c}}_N$ be the conditional law of $(P,\Sigma,A)$ given $\mathfrak T^{\mathsf{pre-c}}$, and let $\mathbb P^{\mathsf{pre-c}}_\star$ be the corresponding marginal of the ideal law of \autoref{not:random-variables}. Then
\[
  \|\widehat{\mathbb P}^{\mathsf{pre-c}}_N-\mathbb P^{\mathsf{pre-c}}_\star\|_{\rm TV}\longrightarrow 0
  \quad\textnormal{almost surely},
  \qquad
  \sqrt N(\widehat{\mathbb P}^{\mathsf{pre-c}}_N-\mathbb P^{\mathsf{pre-c}}_\star)
  \Longrightarrow\mathcal N(0,\Gamma_{\mathsf{pre-c}})
\]
for a covariance matrix $\Gamma_{\mathsf{pre-c}}$. Moreover, there are game-dependent constants $c,C>0$ such that, for every $0<\varepsilon<1$,
\[
  \mathbb P\!\left(\|\widehat{\mathbb P}^{\mathsf{pre-c}}_N-\mathbb P^{\mathsf{pre-c}}_\star\|_{\rm TV}>\varepsilon\right)
  \leq C e^{-cN\varepsilon^2}.
\]
\end{lemma}

\paragraph{Theory in the Ideal Setup.}
Essentially, we modify the expected value computations as
\begin{equation}[rCl]
  \widehat{Q}
  &=&
  \mathbb{E}^\pi(R_{p}(x, \tau) \mid x_p, \sigma a)
  \nonumber\\
  &=&
  \mathbb{E}^\pi(\mathbb{E}^\pi(R_{p}(x, \tau) \mid x_p, \tau^{\mathsf{pre-c}}) \mid x_p, \sigma a)
  \label{eq:rao-blackwell-q-identity}
\end{equation}
\begin{equation}[rCl]
  \mathbb{E}(\mathcal{L})
  &=&
  \mathbb{E}(\mathbb{E}^\pi(\mathcal{L}(x, \Sigma A) \mid \Sigma A))
  \label{eq:rao-blackwell-loss-identity}
\end{equation}
which reduces variance by the Law of Total Variance.
For losses that depend nonlinearly on their targets, substituting these conditional means defines a new Rao--Blackwellized surrogate rather than the exact conditional expectation of the original loss; we expect the lower-variance targets to improve training nevertheless.

\paragraph{Primitives.}
For $\lambda \in [0,1]$, $1 \leq i \leq N$, $\sigma a \preccurlyeq \Tau_i$, $p = o(\sigma)$ and $x_p$ compatible with $\Tau_i$, we write
\begin{equation}\label{eq:rao-blackwell-lambda-return}
  \widehat{Q}^\lambda(x_p, \sigma a)
  :=
  \lambda^{|\Tau_i^{\mathsf{pre-c}} \setminus \sigma|_p} \mathbb{E}^{\pi^{\theta_k}}(R_{p}(x, \tau) \mid x_p, \Tau_i^{\mathsf{pre-c}} \preccurlyeq \tau)
  + (1 - \lambda) \sum_{\substack{\sigma \prec \rho \preccurlyeq \Tau_i \\ o(\rho) = p}}
  \lambda^{|\rho \setminus \sigma|_p - 1} V^{\theta_k}(x_p,\rho)
\end{equation}
and $\widehat{A}^\lambda(x_p) = \widehat{Q}^\lambda(x_p,\Sigma A) - V^{\theta_k}(x_p,\Sigma)$
and $\widehat{V}^\lambda(x_p) = \widehat{Q}^\lambda(x_p,\Sigma A)$.

\paragraph{Monte Carlo Methods.}
After experimentation, we found that a sequential importance sampler (SIS) resulted in the fastest convergence of the Monte Carlo estimates.
Other methods such as rejection sampling or MCMC were found to be less efficient.

\begin{algorithm}[H]
\caption{$(\widehat{\mathbb P}^{\mathfrak r},\widehat V)\gets$~\texttt{PreTerminalMC}$(\mathfrak r,\tau^{\mathsf{pre-c}})$}
\label{alg:mc-preterminal}
\DontPrintSemicolon
\tcc{Black-box Monte Carlo completion from public pre-terminal prefixes.}
\KwParam{$\mathfrak r$, reach cache; $\tau^{\mathsf{pre-c}}$, fixed pre-terminal trajectory}
\KwHyper{Monte Carlo budget}
\KwOut{$(\widehat{\mathbb P}^{\mathfrak r},\widehat V)$}
  Use a Monte Carlo simulation to sample $X \sim \mathbb{P}^\mathfrak{r}(\cdot\mid\tau^{\mathsf{pre-c}})$ with
  \[
    \mathbb{P}^\mathfrak{r}(x\mid\tau^{\mathsf{pre-c}})
    = \frac{\mu(x)\prod_{q \in \mathcal P \cup \{\mathsf c\}} r_q(\tau^{\mathsf{pre-c}} \mid x_q)}
           {\sum_{y\in\mathcal X} \mu(y)\prod_{q \in \mathcal P \cup \{\mathsf c\}} r_q(\tau^{\mathsf{pre-c}} \mid y_q)}
  \]\;
  Conditionally on $X$ and $\tau^{\mathsf{pre-c}}$, sample a terminal completion
  $\Tau\sim\mathbb P^\pi(\cdot\mid X,\tau^{\mathsf{pre-c}})$\;
  Store statistical occurrence $(\widehat{\mathbb P}^{\mathfrak r}(x_p))_{\substack{p \in \mathcal{P} \\ x_p \in \mathcal{X}_p}}$ and statistical average
  \[
    \left(\widehat{V}(x_p) \simeq \mathbb{E}(R_p(X,\Tau) \mid X_p = x_p,\tau^{\mathsf{pre-c}})\right)_{\substack{p \in \mathcal{P} \\ x_p \in \mathcal{X}_p}}
  \]
\Return{$(\widehat{\mathbb P}^{\mathfrak r},\widehat V)$}
\end{algorithm}

\begin{algorithm}[H]
\caption{$(\theta_k)_{k\ge0}\gets$~\texttt{SequentialRaoBlackwellPG}$(\theta_0)$}
\label{alg:sequential-rao-blackwell-pg}
\DontPrintSemicolon
\tcc{Sequential rollout policy-gradient template corresponding to the procedure above.}
\KwParam{$\theta_0$, initial weights; $\theta$, working copy of the next iterate}
\KwHyper{$N$, rollout count; $B$, minibatch size; chosen number of gradient-descent steps per iterate}
\KwOut{$(\theta_k)_{k\ge0}$}
\For{$k=0,1,2,\dots$}{
  Sample $(\Tau_i^{\mathsf{pre-c}}, \mathfrak{r}_i)_{i=1}^N \stackrel{\mathsf{iid}}{\sim} \Tau^{\pi^{\theta_k}}$ using $\algcall{alg:sample-preterminal}{SamplePreTerminal}(\pi^{\theta_k},N)$\;
	For $1 \leq i \leq N$, $p \in \mathcal{P}$ and $x_p \in \mathcal{X}_p$, estimate 
  \begin{equation}[rCl]
        \widehat{\mathbb{P}}(x_p \mid \Tau^{\mathsf{pre-c}}_i)
         &\simeq& \mathbb{P}^\pi(x_p \mid \Tau^{\mathsf{pre-c}}_i)
        \nonumber\\
        \widehat{V}_p(x_p, \Tau^{\mathsf{pre-c}}_i) &\simeq& \mathbb{E}^\pi(R_{p}(x, \tau) \mid x_p, \Tau^{\mathsf{pre-c}}_i)
        \label{eq:sequential-preterminal-estimates}
    \end{equation} 
  using $\algcall{alg:mc-preterminal}{PreTerminalMC}(\mathfrak{r_i},\Tau_i^{\mathsf{pre-c}})$\;
  \For{$i=1,\ldots,N$}{
    By backward propagation, for $\sigma a \preccurlyeq \Tau_i$, $p = o(\sigma)$ and $x_p \in \mathcal{X}_p$, compute $\widehat{Q}(x_p, \sigma a)$ using \autoref{eq:rao-blackwell-lambda-return}\;
    Compute $\widehat{V}(x_p, \sigma)$\;
  }
  $\theta\gets\theta_k$\;
  \For{each gradient-descent step}{
    Sample $(I_b, \Sigma_b, A_b, P_b)_{b=1}^B \stackrel{\mathsf{iid}}{\sim} (I, \Sigma, A, P)$ as in \autoref{def:preterminal-rao-blackwell-decision-law}\;
    Deduce $\widehat{A}_b(-)$, $\widehat{Q}_b(-)$, and $\widehat{V}_b(-)$ at $(\Sigma_b,A_b)$\;
    Estimate
    \begin{equation*}
      \widehat{l}(\theta)
      =
      \frac{1}{B}\sum_{b=1}^B \sum_{x_{P_b}}
        \mathcal{L}\left(
        (\pi^\theta, V^\theta,\pi^{\theta_k}, V^{\theta_k})(x_{P_b}, \Sigma_b),
        A_b,
        \widehat{A}_b(x_{P_b}),
        \widehat{Q}_b(x_{P_b}),
        \widehat{V}_b(x_{P_b})
        \right)
        \cdot
        \widehat{\mathbb{P}}(x_{P_b} \mid \Tau_{I_b}^{\mathsf{pre-c}})
    \end{equation*}\;
    Descend $\theta$ along $\nabla_\theta \widehat{l}(\theta)$\;
  }
  store $\theta_{k+1}=\theta$\;
}
\Return $(\theta_k)_{k\geq 0}$
\end{algorithm}

\subsection{Shared-Seat Approximate-Response Probes}
\label{subsec:shared-seat-approximate-response}

This subsection defines the fifth production workflow.  It is deliberately an
\emph{approximate-response probe}, not an exact best-response solver and not a
certificate that a target profile is near equilibrium.

Fix a target strategy profile $\bar\pi$.  Let $\Theta$ be the configured neural
parameter space and let
\begin{equation}\label{eq:shared-response-class}
  \Pi_{\mathrm{shared}}
  :=\left\{(\pi_p^\theta)_{p\in\mathcal P}:\theta\in\Theta\right\}.
\end{equation}
All seats share the same parameter vector and architecture; the acting seat and
public state are model inputs.  Thus this is one seat-conditioned response
class, not $|\mathcal P|$ independently parameterized response classes.

For any behavioral deviation $\nu_p$ by player $p$, define
\begin{equation}\label{eq:unilateral-gain}
  g_p(\nu_p;\bar\pi)
  :=J_p(\nu_p,\bar\pi_{-p})-J_p(\bar\pi),
\end{equation}
and define the exact best-response gain and NashConv using the full behavioral
strategy class $\Pi_p$:
\begin{equation}\label{eq:exact-br-and-nashconv}
  g_p^*(\bar\pi)
  :=\sup_{\nu_p\in\Pi_p}g_p(\nu_p;\bar\pi),
  \qquad
  \operatorname{NashConv}(\bar\pi)
  :=\sum_{p\in\mathcal P}g_p^*(\bar\pi).
\end{equation}
For a shared probe parameter $\theta$, write
\begin{equation}\label{eq:shared-probe-gain-objective}
  g_p(\theta;\bar\pi):=g_p(\pi_p^\theta;\bar\pi),
  \qquad
  G_{\mathrm{shared}}(\theta;\bar\pi)
  :=\frac1{|\mathcal P|}\sum_{p\in\mathcal P}g_p(\theta;\bar\pi).
\end{equation}
Every $\pi_p^\theta$ is a valid unilateral deviation, so for every fixed
$\theta$,
\begin{equation}\label{eq:shared-probe-lower-bound}
  \sum_{p\in\mathcal P}g_p(\theta;\bar\pi)
  \leq \operatorname{NashConv}(\bar\pi).
\end{equation}
The left side may be negative and remains a lower bound.  No maximum over
$\Theta$, no finite optimization run, and no Monte Carlo estimate may be called
an exact best response.

\paragraph{Training seat law and mixed rollout.}
Independently for every logical root $i$, before private-state or public-event
sampling, draw
\begin{equation}\label{eq:probe-seat-law}
  I_i\sim\mathsf{Unif}(\mathcal P)
\end{equation}
from the counter-based logical key assigned to that root.  The draw is
independent of queue order, physical batching, and GPU placement.  At training
iteration $k$, root $i$ uses the mixed profile
\begin{equation}\label{eq:mixed-probe-rollout-profile}
  \rho^{k,i}_p
  :=
  \begin{cases}[l"l]
    \pi_p^{\theta_k},&p=I_i,\\
    \bar\pi_p,&p\neq I_i.
  \end{cases}
\end{equation}
The target $\bar\pi$ is frozen and identified by its checkpoint content hash.
Only decisions owned by $I_i$ enter the retained decision law; opponent
forwards, actions, and reach updates affect successors but never create
training rows.

Let $\mathcal D_i$ be the selected-seat decisions retained from root $i$.
Conditionally on a fixed rollout batch, the learner samples uniformly from the
pooled set $\bigsqcup_i\mathcal D_i$ and applies exactly the sequential-RB PPO
loss: clipped policy surrogate, actor-type value regression with the configured
value clipping, entropy, and optional legal-action reference KL, all integrated
against the selected seat's completion posterior.  Equivalently, its frozen
batch objective is
\begin{equation}\label{eq:shared-probe-empirical-loss}
  \widehat l_{\mathrm{probe},k}(\theta)
  :=
  \frac{1}{\sum_i|\mathcal D_i|}
  \sum_i\sum_{d\in\mathcal D_i}
  \sum_{x_{I_i}\in\mathcal X_{I_i}}
  \widehat{\mathbb P}_d(x_{I_i})
  \mathcal L_{\mathrm{seqRB},d}(\theta,x_{I_i}).
\end{equation}
The root seat law gives every seat equal ex ante representation, but the
normalization by the realized pooled decision count weights the resulting
state occupancies by retained decisions.  Consequently
\autoref{eq:shared-probe-empirical-loss} is the declared optimization surrogate,
not an unbiased estimator of $G_{\mathrm{shared}}$ or its gradient.  PPO
clipping, value regression, entropy, reference KL, finite batches, and repeated
optimization introduce further surrogate effects.  Any other root-versus-row
weighting is a different response-training contract.

Training terminates only at the configured iteration and optimizer-step budget.
If the final parameters are $\widehat\theta_K$, the result is the finite-budget
shared approximate response $\pi^{\widehat\theta_K}$.  Logical trajectories,
completion samples, minibatches, optimizer steps, target hash, and initialization
rule are part of its reportable scientific identity.  There is no convergence,
exploitability, or optimality stopping certificate.

\paragraph{Public fold cutoff.}
An early selected-seat cutoff is permitted only when there is a publicly
measurable stopping set $\mathcal F_p$ and payoff $F_p$ such that
\begin{equation}\label{eq:probe-public-cutoff}
  R_p(x,\tau)=F_p(\sigma)
  \quad\textnormal{for every }\sigma\in\mathcal F_p,
  \ x\textnormal{ compatible with }\sigma,
  \textnormal{ and terminal }\tau\succeq\sigma.
\end{equation}
In equal-stack cash poker, $\mathcal F_p$ is the set of public states immediately
after $p$ folds and
\begin{equation}\label{eq:poker-probe-fold-payoff}
  F_p(\sigma)=-\operatorname{Committed}_p(\sigma).
\end{equation}
The rollout stops at the first such state because later opponents cannot change
$p$'s payoff.  This is a synthetic probe leaf, not necessarily a terminal poker
state.  Its immutable capture must carry the cutoff identity and selected seat
until it has written $F_p$ and the selected posterior/value contract; these
facts may never be reconstructed by reclassifying the underlying game state.
No opponent decision after the cutoff is sampled or retained.

\paragraph{Paired lower-bound evaluation.}
Evaluation freezes both $\widehat\theta_K$ and $\bar\pi$.  Each consecutive
block of $|\mathcal P|^2$ logical roots contains every
$(\textnormal{probe seat},\textnormal{button position})$ pair exactly once.
For each root $j$, run the mixed profile
$(\pi_{I_j}^{\widehat\theta_K},\bar\pi_{-I_j})$ and the baseline $\bar\pi$
with the same root identity and counter-key namespaces.  The root-private draw
is shared; matching logical addresses reuse keys, while genuinely divergent
paths use their own keyed chance/completion draws.  Record
\begin{equation}\label{eq:paired-probe-difference}
  D_j:=R_{I_j}^{\mathrm{mixed},j}-R_{I_j}^{\mathrm{baseline},j}.
\end{equation}
With $M_p=\#\{j:I_j=p\}$, define
\begin{equation}\label{eq:approximate-nashconv-estimator}
  \widehat g_p:=\frac1{M_p}\sum_{j:I_j=p}D_j,
  \qquad
  \widehat{\operatorname{NashConv}}_{\mathrm{probe}}
  :=\sum_{p\in\mathcal P}\widehat g_p
  =|\mathcal P|\frac1{\sum_pM_p}\sum_jD_j.
\end{equation}
The equality on the right requires the mandated balanced seat counts.  For a
fixed probe checkpoint, $\widehat g_p$ estimates
$g_p(\widehat\theta_K;\bar\pi)$, so the sum estimates the lower-bound quantity
in \autoref{eq:shared-probe-lower-bound}.  The paired Student-$t$ interval
quantifies evaluation sampling error only; it does not include response-class
restriction or optimization error.  Reports must therefore use
``shared approximate response'', ``probe gain'', and ``approximate NashConv
lower bound'', never ``exact BR'' or unqualified ``NashConv''.

\subsection{Variance-Guided Branching}\label{subsec:variance-guided-branching}

Rao--Blackwellization integrates potential private states knowing a public terminal trajectory.
It leaves on the table another substantial source of variance: a state $\sigma a$ might end in many different terminal trajectories.
We suggest here to modify sequential rollout sampling into a Monte Carlo Tree Search (MCTS) style of sampling.
This approach is particularly useful in games because typically there is only one possible initial public state, or a small space thereof (we build this into the PBG formalism where different initial public states are different games).
For example, in our leading example of Texas Hold'em cash poker, the initial public state only contains the stack size of each player (and maybe antes, straddles or other house rules): this is a much smaller space than the space of all potential prompts for a coding model.

We further improve MCTS by suggesting a variance-guided branching strategy: we branch more on the nodes with high downstream reward variance.

\paragraph{Branched Rollouts.}
We enrich the rollout into a branching random walk on the public tree.
The design constraint, encouraged by the Rao--Blackwellized pipeline, is that the enrichment must be \emph{public}: any stored importance weight depending on a sampled hidden assignment would contaminate every posterior downstream (Warning~\ref{warn:no-private-tilting}).

\begin{algorithm}[H]
\caption{$(\mathfrak S^{\beta,\pi},\ell^{\beta,\pi},N,\widetilde N,\mathfrak r)\gets$~\texttt{SampleTree}$(\pi,\beta;C,M_{\tt safe})$}
\label{alg:sample-tree}
\DontPrintSemicolon
\tcc{Sampling procedure for the rollout tree of \cref{def:branched-rollout-tree}.}
\KwParam{$\pi$, rollout profile; $\beta:\mathcal S\setminus\mathcal T\to[1,\infty)$, intensity; optional frozen cost head $C$ and budget $M_{\tt safe}$; optional lower tree-size safeguard $M_{\min}$}
\KwOut{sampled tree $\mathfrak{S}^{\beta,\pi}$, label map $\ell^{\beta,\pi}$, executed and nominal child counts $(N,\widetilde N)$, and node-specific reach caches $\mathfrak r$}
Set $\mathfrak{S}^{\beta,\pi}\gets\{\varnothing\}$, $\ell^{\beta,\pi}(\varnothing)\gets\varnothing$, $\mathfrak r_{\varnothing}\gets(1)_{q\in\mathcal P\cup\{\mathsf c\}}$, and expansion queue $\mathcal F\gets\{\varnothing\}$\;
\While{$\mathcal F\neq\varnothing$}{
  Compute the next forward batch $\mathcal B\subseteq\mathcal F$ and maintain $R_{\mathcal F}=\sum_{\mathsf s\in\mathcal F}C(\ell(\mathsf s))$ incrementally (or $R_{\mathcal F}=0$ if $M_{\tt safe}=\infty$)\;
  \For{$\mathsf s\in\mathcal B$}{
  Set $\sigma\gets\ell^{\beta,\pi}(\mathsf{s})$\;
  \eIf{$\sigma$ is pre-terminal}{
    Set $\mathfrak{E}_{\mathsf{s}}\gets\varnothing$ and $(N(\mathsf s),\widetilde N(\mathsf s))\gets(0,0)$\;
  }{
    Set $q\gets o(\sigma)$\;
    Let $\widetilde N(\mathsf{s}) \sim \lfloor\beta(\sigma)\rfloor + \operatorname{Bernoulli}\big(\beta(\sigma) - \lfloor\beta(\sigma)\rfloor\big)$, and set $N(\mathsf s)\gets1$ if $q\in\mathcal P$ and $|\mathfrak S^{\beta,\pi}|+R_{\mathcal F}>M_{\tt safe}$, and $N(\mathsf s)\gets\widetilde N(\mathsf s)$ otherwise\;
    Set $\mathfrak{E}_{\mathsf{s}}\gets\{\mathsf{a}_1,\ldots,\mathsf{a}_{N(\mathsf{s})}\}$ to be the set of children of $\mathsf{s}$\;
    \For{$\mathsf a\in\mathfrak E_{\mathsf s}$, independently}{
      Sample $X_{\mathsf s,\mathsf a}$ from the posterior represented by the current reach cache,
      \[
        \mathbb P^{\mathfrak r_{\mathsf s}}(x\mid\sigma)
        =
        \frac{\mu(x)\prod_{u\in\mathcal P\cup\{\mathsf c\}}r_{\mathsf s,u}(\sigma\mid x_u)}
             {\sum_{y\in\mathcal X}\mu(y)\prod_{u\in\mathcal P\cup\{\mathsf c\}}r_{\mathsf s,u}(\sigma\mid y_u)}
      \]\;
      Sample
      \[
        E_{\mathsf a}\sim
        \begin{cases}[l"l]
          \pi(\cdot\mid X_{\mathsf s,\mathsf a,q},\sigma) & \textnormal{if }q\in\mathcal P,\\
          \chi(\cdot\mid X_{\mathsf s,\mathsf a},\sigma) & \textnormal{otherwise}
        \end{cases}
      \]\;
      Set $\ell^{\beta,\pi}(\mathsf s\mathsf a)\gets\sigma E_{\mathsf a}$ and $\mathfrak r_{\mathsf s\mathsf a}\gets\mathfrak r_{\mathsf s}$\;
      Set
      \[
        \mathfrak r_{\mathsf s\mathsf a,q}\gets
        \mathfrak r_{\mathsf s\mathsf a,q}\odot
        \begin{cases}[l"l]
          \pi(E_{\mathsf a}\mid-,\sigma) & \textnormal{if }q\in\mathcal P,\\
          \chi(E_{\mathsf a}\mid-,\sigma) & \textnormal{otherwise}
        \end{cases}
      \]\;
      Normalize $\mathfrak r_{\mathsf s\mathsf a,q}$\;
    }
  }
  }
  Set $\mathcal A_{\mathcal B}\gets
  \{\mathsf s\mathsf a:\mathsf s\in\mathcal B,\mathsf a\in\mathfrak E_{\mathsf s}\}$,
  $\mathfrak S^{\beta,\pi}\gets\mathfrak S^{\beta,\pi}\sqcup\mathcal A_{\mathcal B}$,
  $\mathcal F\gets(\mathcal F\setminus\mathcal B)\sqcup\mathcal A_{\mathcal B}$,
  and update $R_{\mathcal F}$ by the corresponding removals and insertions\;
}
Let $N_0(\varnothing)$ be the executed root degree and $\mathfrak S_0$ the
completed tree before lower-size safeguarding.  While
$|\mathfrak S|<M_{\min}$, append one fresh independently sampled root child,
using the next root-child logical ordinal, and expand its complete descendant
subtree by the ordinary $\beta$ rule.  Let $N_+(\varnothing)$ be the final root
degree.  Set every final root-edge transition mass to
$1/N_+(\varnothing)$ and recompute descendant path masses.  Retain
$N_0(\varnothing)$ and $|\mathfrak S_0|$ as protected-cohort metadata; neither
replaces the nominal randomized-rounded count $\widetilde N(\varnothing)$.
\Return{$(\mathfrak{S}^{\beta,\pi},\ell^{\beta,\pi},N,\widetilde N,\mathfrak r)$}
\end{algorithm}

\begin{remark}[Practical Implementation]
  Note that practically speaking, to sample each outgoing event $E_{\mathsf{a}}$, we first sample a temporary private state $x$ conditioned on $\sigma$ (for example, in poker, we do this by rejection sampling), then we sample according to the chance rules and player strategies.
  Since $\beta\geq1$, randomized rounding always gives at least one child; this guaranteed child is a \emph{spine} which preserves every lineage.
  Every abstract sampled child is retained as a distinct tree node, even when two event labels $E_{\mathsf a}$ agree.
  Such nodes may share immutable public-state payloads, but never topology, descendants, random draws, or Monte Carlo completion samples; each outgoing edge has transition probability $1/N(\mathsf s)$.
  At $\beta \equiv 1$ the tree is the set of nodes in a random trajectory $\Tau \sim \mathbb{P}^\pi(\cdot)$.
\end{remark}

\begin{definition}[Rollout tree]\label{def:branched-rollout-tree}
Let $\beta : \mathcal S \setminus \mathcal T \to [1, \infty)$ be an intensity function and $(\mathfrak{S}^{\beta,\pi}, \ell^{\beta,\pi}:\mathfrak{S}^{\beta,\pi}\to\mathcal S, N^{\beta,\pi}:\mathfrak{S}^{\beta,\pi}\to\{0,1,\ldots\},\widetilde N^{\beta,\pi},\mathfrak r) \gets \algcall{alg:sample-tree}{SampleTree}(\pi, \beta;\bot,\infty)$ be a random labeled tree, seen as a Markov chain (with uniform transition probabilities), together with its auxiliary node-specific reach caches.  Here $N^{\beta,\pi}=\widetilde N^{\beta,\pi}$.
We have a natural probability distribution on $\mathcal{X} \times \mathcal{T}$ induced by first sampling a leaf of this Markov chain and then completing its pre-terminal public label:
\begin{equation}\label{eq:branched-terminal-law}
  \mathbb{P}^{\beta, \pi}(x,\tau \mid \mathfrak{S}^{\beta, \pi})
  = \sum_{\ell(\mathsf{t})=\tau^{\mathsf{pre-c}}}
    \left(\prod_{\mathsf{s} \prec \mathsf{t}} \frac{1}{N^{\beta,\pi}(\mathsf{s})}\right)
    \mathbb P^\pi(x,\tau\mid\tau^{\mathsf{pre-c}}).
\end{equation}
From it, we can firstly define $P^{\beta, \pi} \in \mathcal{P}$ by
\begin{equation}\label{eq:branched-player-law}
  \mathbb{P}(P^{\beta, \pi} = p \mid \mathfrak{S}^{\beta, \pi})
  = \frac{\mathbb{E}^{\beta, \pi}(|\tau|_p \mid \mathfrak{S}^{\beta, \pi})}
    {\mathbb{E}^{\beta, \pi}(|\tau|_{\textnormal{player}} \mid \mathfrak{S}^{\beta, \pi})}
\end{equation}
Secondly, we denote $\mathsf{S}^{\beta, \pi}$ for a random node sampled from the stationary distribution of the Markov chain defined on the sub-tree of $\mathfrak{S}^{\beta, \pi}$ on these nodes labeled by a game node owned by $P^{\beta, \pi}$, and we let $\Sigma^{\beta, \pi} = \ell(\mathsf{S}^{\beta, \pi})$ be its label.
Thirdly, we let $\mathsf{A}^{\beta, \pi}$ be a uniform choice in $\mathfrak{E}_{\mathsf{S}^{\beta,\pi}}$ and let $A^{\beta, \pi}$ be its public event label, so $\ell(\mathsf{S}^{\beta,\pi}\mathsf{A}^{\beta,\pi})=\Sigma^{\beta,\pi}A^{\beta,\pi}$.
It's easy to see that the law of $\mathsf{S}^{\beta, \pi}$ is given by
\begin{equation}\label{eq:branched-stationary-node-law}
  \mathbb{P}(\mathsf{S}^{\beta,\pi} = \mathsf{s} \mid P^{\beta, \pi}, \mathfrak{S}^{\beta, \pi})
  = \frac{\one_{o(\ell(\mathsf{s})) = P^{\beta, \pi}}}
         {\mathbb{E}^{\beta, \pi}(|\tau|_{P^{\beta, \pi}} \mid \mathfrak{S}^{\beta, \pi}) \prod_{\mathsf{r} \prec \mathsf{s}} N^{\beta,\pi}(\mathsf{r})}
\end{equation}
which is indeed a probability distribution.
\end{definition}

\begin{lemma}[Single-tree convergence of the branched decision law]\label{lem:branched-decision-law-convergence}
Fix an intensity profile $\bar\beta:\mathcal S\setminus\mathcal T\to[1,\infty)$ and assume that every terminal trajectory contains at least one player node.
For $K\in\mathbb N$, set $\beta_K:=K\bar\beta$, sample a single rollout tree $\mathfrak S_K:=\mathfrak S^{\beta_K,\pi}$ as in \cref{def:branched-rollout-tree}, and let
\[
  \widehat{\mathbb P}_K
  :=
  \mathcal L\!\left(
    P^{\beta_K,\pi},\Sigma^{\beta_K,\pi},A^{\beta_K,\pi}
    \middle|
    \mathfrak S_K
  \right).
\]
Let $\mathbb P_\star^{\mathrm{dec}}:=\mathcal L(P^\pi,\Sigma^\pi,A^\pi)$ be the corresponding marginal of the ideal law in \autoref{not:random-variables}.
Then $|\mathfrak S_K|\to\infty$ and
\[
  \bigl\|\widehat{\mathbb P}_K-\mathbb P_\star^{\mathrm{dec}}\bigr\|_{\mathrm{TV}}
  \xrightarrow[K\to\infty]{\mathbb P}
  0.
\]
\end{lemma}

\begin{lemma}[Terminal node distribution]\label{lem:branched-terminal-distribution}
    We have
    \begin{equation}\label{eq:branched-terminal-distribution}
        \mathbb{E}(\mathbb{P}^{\beta, \pi}(\cdot \mid \mathfrak{S}^{\beta, \pi})) = \mathbb{P}^\pi(\cdot)
    \end{equation}
\end{lemma}

\begin{lemma}[Branch independence]\label{lem:branched-branch-independence}
    Fix a sampled-tree node $\mathsf{s}$ with public label $\sigma$ and sampled outgoing-action set $\mathfrak{E}_{\mathsf{s}}$, and put $m=N(\mathsf{s})=|\mathfrak{E}_{\mathsf{s}}|$.
    Conditional on $\sigma$ and $m$, the rooted child branches
    \[
        \left(\ell(\mathsf{s}\mathsf{a}),\mathfrak{S}_{\mathsf{s}\mathsf{a}\preccurlyeq}^{\beta,\pi}\right)_{\mathsf{a}\in\mathfrak{E}_{\mathsf{s}}}
    \]
    are mutually independent; in particular, any statistic of one child branch is independent of any statistic of the sibling forest, conditional on the public parent and the chosen child label.
\end{lemma}

\begin{warning}[Don't sample a true hand!]\label{warn:no-private-tilting}
    One might consider sampling $X$ first and then generating the sub-tree using the knowledge that $X$ is the true private state.
    While this still satisfies \autoref{lem:branched-terminal-distribution} marginally, it violates \autoref{lem:branched-branch-independence}.
    The sub-trees sampled by this process would be ``self-correlated'': for $\mathsf{a}\neq\mathsf{b}\in \mathfrak{E}_{\mathsf{s}}$, the sub-trees $\mathfrak{S}_{\mathsf{s}\mathsf{a} \preccurlyeq}$ and $\mathfrak{S}_{\mathsf{s}\mathsf{b} \preccurlyeq}$ share the same sampled private state and would not be independent.
    In particular, the sibling average used as the leave-one-out baseline for $\widehat{A}^\lambda$ would no longer be independent of the target branch after conditioning on the target action label.
\end{warning}

\paragraph{Queued execution laws that supersede the independent-tree estimator.}
The queued rewrite also exposes two explicitly different branching laws. They
must retain distinct mode identity and must not inherit claims that require
\autoref{lem:branched-branch-independence}.

In sampled-private branching only (a branched rollout with the
\texttt{sampled} private estimator), each logical root first draws one joint
private profile. Every action and chance draw in that root's complete rollout
tree is conditioned on that stored profile. Siblings are therefore
intentionally correlated. This mode is a deliberate negative-control
experiment intended to demonstrate that the estimator is a bad idea. It is a
biased estimator: it does not satisfy
\autoref{lem:branched-branch-independence}, does not use that lemma to justify a
leave-one-out baseline, and must never be reported as the independent-tree law
of \cref{alg:sample-tree}.

In branched Rao--Blackwell execution, a node with $N(\mathsf s)\geq1$ assigns
its inherited posterior profile to logical child zero and draws exactly
$N(\mathsf s)-1$ fresh posterior profiles for the remaining children. Each
public action or chance label is then drawn conditional on its child's
profile, and that profile is propagated to the child. The inherited profile
has the node posterior law by construction; the fresh profiles are
independent draws from the same law. This child-zero reuse is the authoritative
queued sampling law and avoids one rejection sample per expansion.

A run may contain any positive number of logical roots. Root $r$ owns a local
distinguished root $(r,\varnothing)$, intrinsic root fields, and an independent
logical RNG address. Root batching is an execution dimension, not a branching
degree or a new estimator mode. At $\beta\equiv1$, every root has
$N=\widetilde N=N^{\tt tgt}=1$ at every node and is exactly one sequential
trajectory.

\paragraph{Primitive Functions.}
Back to a synchronous rollout training, we fix weights $\theta_0,\ldots, \theta_k$.
We let $\mathfrak{S}$, $\ell$, $P$, $\mathsf{S}$, $\Sigma$, $\mathsf{A}$, $A$, $\mathsf{T}$, $\Tau^{\mathsf{pre-c}}$ and $X$ be as above for $\pi = \pi^{\theta_k}$.
Throughout this paragraph, $\mathbb E^{\beta,\pi}$ is conditional on the realized sampled labeled tree $\mathfrak S$; conditions involving $\mathfrak S$ are displayed only where they clarify the sampled-node law.
For a sampled-tree node $\mathsf{s}$, write $\mathfrak{E}_{\mathsf{s}}$ for its sampled outgoing-action set, so that the children are $\mathsf{s}\mathsf{a}$ for $\mathsf{a}\in\mathfrak{E}_{\mathsf{s}}$.
For a player node $\mathsf{s}\in \mathfrak{S}$ with $\ell(\mathsf{s})=\sigma$, a sampled action $\mathsf{a}\in\mathfrak{E}_{\mathsf{s}}$ with $\ell(\mathsf{s}\mathsf{a})=\sigma a_{\mathsf{a}}$, $p=o(\sigma)$ and $x_p\in\mathcal X_p$, we define the estimators
\begin{equation}[rCl]
    \widehat{Q}^\lambda(x_p, \mathsf{s}\mathsf{a}) &=&
    \mathbb{E}^{\beta, \pi}\left(
        \lambda^{|\tau \setminus \sigma|_p} R_p(x, \tau) \middle| x_p, \mathsf{s}\mathsf{a}, \mathfrak{S}
    \right)\nonumber\\
    &&\quad+
    (1 - \lambda) \mathbb{E}^{\beta, \pi}(|\tau \setminus \sigma|_{p} \mid x_p, \mathsf{s}\mathsf{a}, \mathfrak{S})
    \mathbb{E}\left(
        \lambda^{|\Sigma \setminus \sigma|_p - 1} V^{\theta_k}(x_p,\Sigma)
        \middle|
        \mathsf{s}\mathsf{a} \preccurlyeq \mathsf{S}, P = p, \mathfrak{S}
    \right)
    \nonumber\\
    \widehat{V}^\lambda(x_p, \mathsf{s})
    &=&
    \mathbb{E}^{\beta, \pi}\left(
        \lambda^{|\tau \setminus \sigma|_p} R_p(x, \tau) \middle| x_p, \mathsf{s}, \mathfrak{S}
    \right)\nonumber\\
    &&\quad+
    (1 - \lambda) \mathbb{E}^{\beta, \pi}(|\tau \setminus \sigma|_{p} \mid \mathsf{s}, \mathfrak{S})
    \mathbb{E}\left(
        \lambda^{|\Sigma \setminus \sigma|_p - 1} V^{\theta_k}(x_p,\Sigma)
        \middle|
        \mathsf{s} \prec \mathsf{S}, P = p, \mathfrak{S}
    \right)
    \nonumber\\
    \widehat{V}^{\lambda,\neg\mathsf{a}}(x_p, \mathsf{s})
    &=&
    \mathbb{E}^{\beta, \pi}\left(
        \lambda^{|\tau \setminus \sigma|_p+1} R_p(x, \tau)
        \middle|
        x_p,\mathsf{s} ,\neg \mathsf{s}\mathsf{a}, \mathfrak{S}
    \right)\nonumber\\
    &&\quad+
    (1 - \lambda)
    \mathbb{E}^{\beta, \pi}\left(
        1+ |\tau \setminus \sigma|_{p}
        \middle|
        x_p, \mathsf{s}, \neg \mathsf{s}\mathsf{a},\mathfrak{S}
    \right)
    \mathbb{E}\left(
        \lambda^{|\Sigma \setminus \sigma|_p} V^{\theta_k}(x_p,\Sigma)
        \middle|
        \mathsf{s} \preccurlyeq \mathsf{S},\mathsf{s}\mathsf{a}\not\preccurlyeq \mathsf{S},P = p,\mathfrak{S}
    \right)
    \quad\textnormal{if }N(\mathsf{s})>1
    \nonumber\\
    \widehat{V}^{\lambda,\neg\mathsf{a}}(x_p, \mathsf{s})
    &=&
    V^{\theta_k}(x_p,\sigma)
    \quad\textnormal{if }N(\mathsf{s})=1
    \nonumber\\
    \widehat{A}^\lambda(x_p,\mathsf{s}\mathsf{a})
    &=&
    \widehat{Q}^\lambda(x_p,\mathsf{s}\mathsf{a})
    -
    \widehat{V}^{\lambda,\neg\mathsf{a}}(x_p,\mathsf{s})
    \label{eq:branched-targets}
\end{equation}
\begin{lemma}[Unnormalised bucket moments]\label{lem:branched-bucket-moments}
  All of the above estimators are unbiased when multiplied by $\one_{x_p}$.
\end{lemma}
For $N(\mathsf{s})>1$, \autoref{lem:branched-branch-independence} makes the sibling average entering $\widehat{V}^{\lambda,\neg\mathsf{a}}$ independent of the target branch after conditioning on the target action label, which is needed for the lemma above.
All of these targets can be computed by backward propagation (see algorithm below).

For a sampled $(P, \mathsf{S}, \Sigma)$, a sampled action $\mathsf{a}\in\mathfrak{E}_{\mathsf{S}}$ labeled by $\Sigma a_{\mathsf{a}}$, and a private state $x_P \in \mathcal{X}_P$, we use the local aliases $\widehat{Q}(x_P, \mathsf{a}) := \widehat{Q}^\lambda(x_P, \mathsf{S}\mathsf{a})$, $\widehat{A}(x_P, \mathsf{a}) := \widehat{A}^\lambda(x_P, \mathsf{S}\mathsf{a})$ and $\widehat{V}(x_P) := \widehat{V}^\lambda(x_P, \mathsf{S})$.

\paragraph{Variance-Based Intensities using KKT.}
Ideally, we'd like to pick $\beta$ so to reduce the variance of the gradient estimator.
In practice, this is too noisy to be useful, so we use the following proxy.
Firstly, we assume that the GAE horizon is $\lambda = 1$ for the variance calculation, so as to simplify notation: in practice this means that we are controlling the Monte Carlo source of variance in the estimators.
Secondly, we define the vector Monte Carlo value estimator
\[
  \widehat{\mathbf{V}}(\mathsf{s})
  =
  \mathbb{E}^{\beta, \pi}\left(
    \mathbf{R}(x, \tau) \middle| \mathsf{s}, \mathfrak{S}
  \right)
\]
By Cauchy-Schwarz, we can bound the error of the gradient estimator by the squared error $\operatorname*{SqEr}(\mathbf V \mid \mathfrak{S}) = \mathbb{E}\left( \|\widehat{\mathbf V}(\mathsf{S}) - \mathbb{E}(\mathbf{R} \mid \Sigma)\|_2^2 \mid \mathfrak{S}\right)$.
A reasonnable proxy is then to minimize the expected squarred error:
\begin{equation*}[rCl]
  \mathbb{E}\left(
    \mathbb{E}(|\tau|_{\textnormal{player}} \mid \mathfrak{S})\operatorname*{SqEr}(\widehat{\mathbf{V}} \mid \mathfrak{S})
  \right)
  &=&
  \mathbb{E}^\pi(|\tau|_{\textnormal{player}}) \mathbb{E}(\operatorname*{trCov}(\widehat{\mathbf{V}} \mid \Sigma))
  \\
  &\simeq&
  \sum_{\sigma\in\mathcal S\setminus\mathcal T}
  \mathbb{P}^\pi(\sigma \preccurlyeq \tau) \mathbf{U}(\sigma)
  \left(
    \sum_{\substack{\rho \preccurlyeq \sigma \\ o(\rho)\in\mathcal P}}
    \frac{1}{\prod_{\rho \preccurlyeq \alpha \preccurlyeq \sigma} \beta(\alpha)}
  \right)
\end{equation*}
for $\mathbf{U}(\sigma) = \operatorname*{trCov}(\mathbb{E}(\mathbf{R} \mid \sigma A) \mid \sigma)$, the portion of the variance coming from the next decision only,
subject to controlling the expected size of the sampled tree:
\[
  \mathbb{E}(|\mathfrak{S}|)
  =
  1+
  \sum_{\sigma\in\mathcal S\setminus\mathcal T}
  \mathbb{P}^\pi(\sigma \preccurlyeq \tau)
  \prod_{\rho \preccurlyeq \sigma} \beta(\rho)
  \leq M
\]
for a budget $M$.

We see this as a \texttt{KKT} fixed point problem in the player-node variables
$(\beta(\sigma))_{o(\sigma)\in\mathcal P}$, keeping chance-node intensities fixed,
with rescaled inverse Lagrange multiplier $\kappa$.
Throughout this calculation, $\beta(\sigma)>0$ is a deterministic continuous
mean-field approximation to the randomized offspring count $N(\mathsf s)$ at
$\ell(\mathsf s)=\sigma$.
Within this mean-field approximation, we compute the derivatives
\begin{equation*}[rCl]
  \frac{\partial}{\partial (1/\beta(\sigma))}
  \left(
      \mathbb{E}(\operatorname*{trCov}(\widehat{\mathbf V} \mid \Sigma))
  \right)
  &=&
  \frac{\mathbb{P}^\pi(\sigma \preccurlyeq \tau)}{\mathbb{E}^\pi (|\tau|_{\textnormal{player}})}
    \left(
      \sum_{\substack{\rho \preccurlyeq \sigma \\ o(\rho)\in\mathcal P}}
      \frac{1}{\prod_{\rho \preccurlyeq \alpha \prec \sigma} \beta(\alpha)}
    \right)
    \left(
      \sum_{\sigma \preccurlyeq \eta}
      \frac
      {\mathbf{U}(\eta) \mathbb{P}^\pi(\eta \preccurlyeq \tau \mid \sigma)}
      {\prod_{\sigma \prec \alpha \preccurlyeq \eta} \beta(\alpha)}
    \right)
    \\
    &=&
    \frac{\mathbb{P}^\pi(\sigma \preccurlyeq \tau)}{\mathbb{E}^\pi (|\tau|_{\textnormal{player}})}
    \left(
      \sum_{\substack{\rho \preccurlyeq \sigma \\ o(\rho)\in\mathcal P}}
      \frac{1}{\prod_{\rho \preccurlyeq \alpha \prec \sigma} \beta(\alpha)}
    \right)
    \operatorname*{trCov}(\widehat{\mathbf V} \mid \sigma, N(\mathsf s) = 1)
  \\
    \frac{\partial}{\partial \beta(\sigma)}
  \left(
      \mathbb{E}(|\mathfrak{S}|)
  \right)
  &=&
  \left(
    \mathbb{P}^\pi(\sigma \preccurlyeq \tau)
    \prod_{\rho \prec \sigma} \beta(\rho)
  \right)
  \mathbb{E}(|\mathfrak{S}_{\mathsf{s}\preccurlyeq}|-1
  \mid \ell(\mathsf{s})=\sigma,N(\mathsf{s})=1)
\end{equation*}
This motivates us to define two new positive heads whose regression estimands
are the arithmetic conditional means of the targets defined below:
\begin{enumerate}
  \item \textbf{Trace-Variance Head:}
    $T^\theta(\sigma) \simeq
    \mathbb E(\widehat T(\mathsf s)\mid
    \ell(\mathsf s)=\sigma,N(\mathsf s)\geq2)$, for which
    $T^\theta(\sigma)=\operatorname{sigmoid}(t^\theta(\sigma))$ is appropriate;
  \item \textbf{Cost Head:}
    $C^\theta(\sigma) \simeq
    \mathbb E(\widehat C^1(\mathsf s)\mid
    \ell(\mathsf s)=\sigma)$, for which
    $C^\theta(\sigma)=\exp(c^\theta(\sigma))$ is appropriate.
\end{enumerate}
Then the continuous KKT update is
\begin{equation}
  \beta_k(\sigma) \gets \clip\left(\sqrt{\kappa_k
    \frac{T^{\theta_k}(\sigma)}{C^{\theta_k}(\sigma)}
    \frac{\left(
      \sum_{\substack{\rho \preccurlyeq \sigma \\ o(\rho)\in\mathcal P}}
      \frac{1}{\prod_{\rho \preccurlyeq \alpha \prec \sigma} \beta_k(\alpha)}
    \right)}
    {\prod_{\rho \prec \sigma} \beta_k(\rho)}},1,\beta_{\max}\right)
\end{equation}
for $o(\sigma)\in\mathcal P$, while chance-node intensities remain fixed.

\paragraph{Queue budget safeguard.}
\Cref{alg:sample-tree} stores the nominal rounded count $\widetilde N$ but
executes $N=1$ for one player-node forward batch whenever the one-child queue
forecast exceeds $M_{\tt safe}$.  This is a pressure safeguard rather than a
hard memory bound.  The $T$ and $C$ estimands remain those of the nominal tree;
\cref{alg:branched-value} corrects them using $(N,\widetilde N)$.
For $0<\gamma_\kappa,\eta_\kappa\leq1$, let $\overline M_k$ be an exponential
moving average of the counterfactual size $M_k^{\tt cf}$ defined in
\cref{alg:branched-rao-blackwell-pg}.  After each adaptive rollout, update
\begin{equation}[rCl]
  \overline M_k
  &\gets&
  (1-\gamma_\kappa)\overline M_{k-1}+\gamma_\kappa M_k^{\tt cf}
  \\
  \log\kappa_{k+1}
  &\gets&
  \clip\left(
    \log\kappa_k+\eta_\kappa\log\frac{M}{\overline M_k},
    \log\kappa_{\min},\log\kappa_{\max}
  \right).
\end{equation}
The controller evaluates the KKT update in log space.  Randomized rounding
satisfies $\mathbb E(\widetilde N(\mathsf s)\mid\ell(\mathsf s)=\sigma)=\beta(\sigma)$;
the mean-field KKT ignores the corresponding reciprocal-moment correction.
Practically speaking, we can't immediately use these heads in training before they are somewhat trained.
So we warmup the training on a fixed intensity $\beta_{\tt fix}$ before starting this iterative process.

The lower tree-size safeguard is an outcome-dependent stopping rule and is
therefore different from the queue-pressure correction above. At the root,
all learning targets are reduced over only the first
$N_0(\varnothing)$ children. In particular, writing
$\mathfrak E^0_{\varnothing}$ for that protected prefix,
\[
  \widehat C^\lambda(\varnothing)
  =\frac1{N_0(\varnothing)}
    \sum_{\mathsf a\in\mathfrak E^0_{\varnothing}}
    \left(1+\widetilde N(\mathsf a)\widehat C^\lambda(\mathsf a)\right).
\]
The root trace sample covariance and its average child $\Delta$ correction are
likewise computed only over $\mathfrak E^0_{\varnothing}$, and the trace target
is absent when $N_0(\varnothing)=1$, regardless of the final augmented degree.
The same protected prefix is used for the root posterior, value, and
leave-one-out advantage targets. Refill branches remain ordinary learning data
strictly below the root. This prevents the size stopping rule from changing
root estimands when subtree size is correlated with value or cost.
For the target equations below, define
$\mathfrak E^{\tt tgt}_{\mathsf s}=\mathfrak E^0_{\varnothing}$ at a refilled
root and $\mathfrak E^{\tt tgt}_{\mathsf s}=\mathfrak E_{\mathsf s}$ otherwise,
and let $N^{\tt tgt}(\mathsf s)=|\mathfrak E^{\tt tgt}_{\mathsf s}|$.

To train the trace-variance head, for every sampled node $\mathsf{s}$ labeled by
$\sigma$ with $N^{\tt tgt}(\mathsf{s})\geq 2$, first compute the executed-tree trace
sample covariance
\begin{equation}
  \widehat{T}^{\tt exec}(\mathsf{s})
  =
    \frac{1}{N^{\tt tgt}(\mathsf{s})-1}
  \sum_{\mathsf{a}\in\mathfrak{E}^{\tt tgt}_{\mathsf{s}}}
  \left\|
    \widehat{\mathbf V}(\mathsf{s}\mathsf{a})
    -\widehat{\mathbf V}(\mathsf{s})
  \right\|_2^2
  \label{eq:branched-trace-variance-target}
\end{equation}
and set it to $\bot$ when $N^{\tt tgt}(\mathsf s)=1$.  Using the same GAE horizon
$\lambda\in[0,1]$ as above, $\widehat C^\lambda(\mathsf s)$ estimates nominal
per-child subtree size and $\widehat C^1$ is its Monte Carlo target.  The
counterfactual recurrences are given directly in \cref{alg:branched-value};
the $T$ correction is first-order when the shared queue couples later sibling
work.
For a given sample, we let $\widehat T = \widehat T(\Sigma)$ and
$\widehat C^\lambda = \widehat C^\lambda(\Sigma)$ and introduce the
Gamma-deviance losses
\begin{equation}[rCl]
  \mathcal L_T[\widehat T]
  &=&
  \begin{cases}[l"l]
    c_T\left(
      \dfrac{\widehat T+\varepsilon_T}{T^\theta(\Sigma)+\varepsilon_T}
      +\log\dfrac{T^\theta(\Sigma)+\varepsilon_T}{\widehat T+\varepsilon_T}-1
    \right) & \textnormal{if $N^{\tt tgt}(\mathsf S)\geq2$},\\
    0 & \textnormal{if $N^{\tt tgt}(\mathsf S)=1$},
  \end{cases}
  \\
  \mathcal L_C[\widehat C^\lambda]
  &=&
  c_C\left(
    \dfrac{\widehat C^\lambda}{C^\theta(\Sigma)}
    +\log\dfrac{C^\theta(\Sigma)}{\widehat C^\lambda}-1
  \right).
  \label{eq:branched-intensity-head-losses}
\end{equation}
Their expected optima are the arithmetic conditional means required by the
KKT derivatives, while their relative scale makes very small and very large
targets have comparable influence.

\begin{algorithm}[H]
\caption{$(\widehat{\mathbb P}^{\mathfrak r},\widehat V,\widehat{\mathbf V})\gets$~\texttt{PreTerminalMomentMC}$(\mathfrak r,\tau^{\mathsf{pre-c}})$}
\label{alg:mc-preterminal-moments}
\DontPrintSemicolon
\tcc{Black-box Monte Carlo completion from public pre-terminal prefixes.}
\KwParam{$\mathfrak r$, leaf-specific reach cache; $\tau^{\mathsf{pre-c}}$, fixed pre-terminal trajectory}
\KwHyper{Monte Carlo budget}
\KwOut{$(\widehat{\mathbb P}^{\mathfrak r},\widehat V,\widehat{\mathbf V})$}
  Use a Monte Carlo simulation to sample $X\sim\mathbb P^{\mathfrak r}(\cdot\mid\tau^{\mathsf{pre-c}})$ with
  \[
    \mathbb P^{\mathfrak r}(x\mid\tau^{\mathsf{pre-c}})
    =
    \frac{\mu(x)\prod_{q\in\mathcal P\cup\{\mathsf c\}}r_q(\tau^{\mathsf{pre-c}}\mid x_q)}
         {\sum_{y\in\mathcal X}\mu(y)\prod_{q\in\mathcal P\cup\{\mathsf c\}}r_q(\tau^{\mathsf{pre-c}}\mid y_q)}
  \]\;
  Conditionally on $X$ and $\tau^{\mathsf{pre-c}}$, sample $\Tau\sim\mathbb P^\pi(\cdot\mid X,\tau^{\mathsf{pre-c}})$\;
  Store statistical occurrence $(\widehat{\mathbb P}^{\mathfrak r}(x_p))_{\substack{p \in \mathcal{P} \\ x_p \in \mathcal{X}_p}}$, statistical average
  \[
    \left(\widehat{V}(x_p) \simeq \mathbb{E}(R_p(X,\Tau) \mid X_p = x_p)\right)_{\substack{p \in \mathcal{P} \\ x_p \in \mathcal{X}_p}}
  \]\;
  Compute $\widehat{\mathbf{V}} \gets \left(\widehat{\mathbb P}^{\mathfrak r}(-_p) \cdot \widehat{V}(-_p)\right)_{p \in \mathcal{P}}$\;
\Return{$(\widehat{\mathbb P}^{\mathfrak r},\widehat V,\widehat{\mathbf V})$}
\end{algorithm}

In the backward implementation, tuple operations are componentwise.
For any tuple-valued child expression $F(\bullet)$, write
\[
  \operatorname{Avg}_{\mathsf{s}}[F(\bullet)]
  :=
  \frac{1}{N^{\tt tgt}(\mathsf{s})}
  \sum_{\mathsf{a}\in\mathfrak{E}^{\tt tgt}_{\mathsf{s}}}F(\mathsf{s}\mathsf{a})
\]
and
\[
  \operatorname{Avg}_{\mathsf{s},\neg\mathsf{a}}[F(\bullet)]
  :=
  \begin{cases}[l"l]
    \displaystyle
    \frac{1}{N^{\tt tgt}(\mathsf{s})-1}
    \sum_{\mathsf{b}\in\mathfrak{E}^{\tt tgt}_{\mathsf{s}}\setminus\{\mathsf{a}\}}F(\mathsf{s}\mathsf{b})
    & \textnormal{if }N^{\tt tgt}(\mathsf{s})>1,\\
    \bot & \textnormal{if }N^{\tt tgt}(\mathsf{s})=1.
  \end{cases}
\]
For tuples $z,h$ of the same shape, set
\[
  \operatorname{Mix}_{\lambda}(z,h)
  :=
  \begin{cases}[l"l]
    \lambda z+(1-\lambda)h & \textnormal{if }z\textnormal{ is present},\\
    h & \textnormal{if }z=\bot.
  \end{cases}
\]
The notation $-_p$ denotes the full private-state vector indexed by $x_p\in\mathcal X_p$.
For posterior mass $w$ and posterior-weighted moment $m$, write
\[
  \operatorname{Cond}(m,w):=
  \begin{cases}[l"l]
    m/w & \textnormal{if }w>0,\\
    0 & \textnormal{if }w=0,\\
    \bot & \textnormal{if }m=w=\bot.
  \end{cases}
\]
In numerical implementations the denominator may be clamped to a small positive constant.

\begin{algorithm}[H]
\caption{$\mathcal O \gets$~\texttt{BranchedValue}$(\mathfrak{S}, V, T, C, \widetilde N, \mathfrak r)$}
\label{alg:branched-value}
\DontPrintSemicolon
\tcc{Backward Propagation of Value Estimators.}
\KwParam{$\mathfrak{S}$, executed labeled tree with child counts $N$; $V$, value head; $T$, trace-variance head; $C$, cost head; $\widetilde N$, nominal child counts; $\mathfrak r$, node-specific reach caches}
\KwHyper{$\lambda$, GAE horizon}
\KwOut{target table $\mathcal O$}
Initialize target table $\mathcal O\gets\varnothing$\;
\For{terminal $\mathsf{t} \in \mathfrak{S}$ labeled by pre-terminal $\tau^{\mathsf{pre-c}}$}{
    Monte Carlo estimation of
    \begin{equation}[rCl]
        \widehat{\mathbb{P}}(x_p \mid \mathsf{t}) &\simeq& \mathbb{P}^\pi(x_p \mid \tau^{\mathsf{pre-c}})
        \nonumber\\
        \widehat{V}_p(x_p, \mathsf{t}) &\simeq& \mathbb{E}^\pi(R_{p}(x, \tau) \mid \tau^{\mathsf{pre-c}}, x_p)
        \nonumber\\
        \widehat{\mathbf V}(\mathsf{t}) &\simeq& \mathbb{E}^\pi(\mathbf{R}(x,\tau)\mid \tau^{\mathsf{pre-c}})
        \label{eq:branched-moment-terminal-estimates}
    \end{equation}   
    for $p \in \mathcal{P}$ and $x_p \in \mathcal{X}_p$ using $\algcall{alg:mc-preterminal-moments}{PreTerminalMomentMC}(\mathfrak r_{\mathsf t},\tau^{\mathsf{pre-c}})$\;
    Set $\widehat{M}_p(-_p,\mathsf t)\gets
    \widehat{\mathbb P}(-_p\mid\mathsf t)\widehat V_p(-_p,\mathsf t)$ for every $p\in\mathcal P$\;
    Set $(\widehat C^\lambda(\mathsf t),\Delta(\mathsf t))\gets(0,0)$\;
    Store $\mathcal O[\mathsf t].\mathrm{posterior}\gets
    (\widehat{\mathbb P}(-_p\mid\mathsf t))_{p\in\mathcal P}$\;
}
\For{node $\mathsf{s} \in \mathfrak{S}$ non-terminal, backward}{
  Set $\sigma\gets\ell(\mathsf s)$ and
  $\overline\Delta(\mathsf s)\gets\operatorname{Avg}_{\mathsf s}[\Delta(\bullet)]$\;
  Compute
  \begin{equation}[rCl]
    \left(
      \left(\widehat{\mathbb{P}}(-_p \mid \mathsf{s}),\widehat{M}_p(-_p,\mathsf{s})\right)_{p\in\mathcal P},
      \widehat{\mathbf V}(\mathsf{s}),
      \widehat{C}^\lambda(\mathsf s)
    \right)
    &\gets&
    \operatorname{Avg}_{\mathsf{s}}\left[
      \left(
        \left(\widehat{\mathbb{P}}(-_p \mid \bullet),\widehat{M}_p(-_p,\bullet)\right)_{p\in\mathcal P},
        \widehat{\mathbf V}(\bullet),
        1 + \widetilde N(\mathsf s\bullet)\widehat C^\lambda(\bullet)
      \right)
    \right]
    \nonumber\\
    \label{eq:branched-backward-updates}
  \end{equation}
  Store $\mathcal O[\mathsf s].\mathrm{posterior}\gets
  (\widehat{\mathbb P}(-_p\mid\mathsf s))_{p\in\mathcal P}$\;
  \If{$o(\sigma) \in \mathcal{P}$}{
      Set $\widehat T(\mathsf s)\gets
      \max\{0,\widehat T^{\tt exec}(\mathsf s)-\overline\Delta(\mathsf s)\}$
      if $N^{\tt tgt}(\mathsf s)\geq2$ and to $\bot$ otherwise; set
      $T^*(\mathsf s)\gets\widehat T(\mathsf s)$ if present and to $T(\sigma)$ otherwise; set
      $\Delta(\mathsf s)\gets(T^*(\mathsf s)+\overline\Delta(\mathsf s))/N^{\tt tgt}(\mathsf s)-T^*(\mathsf s)/\widetilde N(\mathsf s)$\;
      $\left(\widehat{Q}^\lambda(-_{o(\sigma)},\mathsf{s}\mathsf{a})\right)_{\mathsf{a}\in\mathfrak{E}^{\tt tgt}_{\mathsf{s}}}
      \gets
      \left(
        \operatorname{Cond}\left(
          \widehat M_{o(\sigma)}(-_{o(\sigma)},\mathsf{s}\mathsf{a}),
          \widehat{\mathbb P}(-_{o(\sigma)}\mid\mathsf{s}\mathsf{a})
        \right)
      \right)_{\mathsf{a}\in\mathfrak{E}^{\tt tgt}_{\mathsf{s}}}$\;
      $\left(\widehat{V}^{\lambda,\neg\mathsf{a}}(-_{o(\sigma)},\mathsf{s})\right)_{\mathsf{a}\in\mathfrak{E}^{\tt tgt}_{\mathsf{s}}}
      \gets
      \left(
        \operatorname{Mix}_{\lambda}\left(
          \operatorname{Cond}\left(
            \operatorname{Avg}_{\mathsf{s},\neg\mathsf{a}}\left[\widehat{M}_{o(\sigma)}(-_{o(\sigma)},\bullet)\right],
            \operatorname{Avg}_{\mathsf{s},\neg\mathsf{a}}\left[\widehat{\mathbb P}(-_{o(\sigma)}\mid\bullet)\right]
          \right),
          V(-_{o(\sigma)},\sigma)
        \right)
      \right)_{\mathsf{a}\in\mathfrak{E}^{\tt tgt}_{\mathsf{s}}}$\;
      $\left(\widehat{A}^\lambda(-_{o(\sigma)},\mathsf{s}\mathsf{a})\right)_{\mathsf{a}\in\mathfrak{E}^{\tt tgt}_{\mathsf{s}}}
      \gets
      \left(\widehat{Q}^\lambda(-_{o(\sigma)},\mathsf{s}\mathsf{a})-\widehat{V}^{\lambda,\neg\mathsf{a}}(-_{o(\sigma)},\mathsf{s})\right)_{\mathsf{a}\in\mathfrak{E}^{\tt tgt}_{\mathsf{s}}}$\;
      $\widehat V_{o(\sigma)}(-_{o(\sigma)},\mathsf s)
      \gets\operatorname{Cond}(\widehat M_{o(\sigma)}(-_{o(\sigma)},\mathsf s),\widehat{\mathbb P}(-_{o(\sigma)}\mid\mathsf s))$\;
      Store current targets $\mathcal O[\mathsf s].\mathrm{targets} \gets \left((\widehat{A}^\lambda(-_{o(\sigma)},\mathsf{s}\mathsf{a}),\widehat{Q}^\lambda(-_{o(\sigma)},\mathsf{s}\mathsf{a}))_{\mathsf{a}\in\mathfrak{E}^{\tt tgt}_{\mathsf{s}}},\widehat V_{o(\sigma)}(-_{o(\sigma)},\mathsf s),\widehat T(\mathsf s),\widehat C^\lambda(\mathsf s)\right)$\;
      Overwrite with the bootstrapped private-state moment propagated to ancestors:
      \begin{equation}[rCl]
        \widehat M_{o(\sigma)}(-_{o(\sigma)},\mathsf s)
        &\gets&
        \operatorname{Mix}_{\lambda}\left(
          \widehat M_{o(\sigma)}(-_{o(\sigma)},\mathsf s),
          \widehat{\mathbb P}(-_{o(\sigma)}\mid\mathsf s)V(-_{o(\sigma)},\sigma)
        \right)
        \\
        \widehat C^\lambda(\mathsf s)
        &\gets&
        \operatorname{Mix}_{\lambda}\left(
          \widehat C^\lambda(\mathsf s),
          C(\sigma)
        \right)
      \end{equation}
  }
  \Else{
    Set $\Delta(\mathsf s)\gets\overline\Delta(\mathsf s)/N^{\tt tgt}(\mathsf s)$\;
  }
}
\Return{target table $\mathcal O$}
\end{algorithm}

Note that returning $\widehat{Q}$ is only necessary for the deferred \texttt{NeuRD} research extension.
The production PPO path returns $\widehat{A}$ and does not store $\widehat Q$ by default.

\begin{algorithm}[H]
\caption{$(\theta_k)_{k\ge0}\gets$~\texttt{BranchedRaoBlackwellPG}$(\theta_0)$}
\label{alg:branched-rao-blackwell-pg}
\DontPrintSemicolon
\tcc{Tree-valued generalization of the sequential Rao--Blackwellized template; $\beta\equiv1$ recovers it exactly.}
\tcc{The implementation exposes fixed and post-warmup dynamic intensity modes; both train actor-type $V$ and public $T/C$ heads with PPO advantage targets.}
\KwParam{$\theta_0$, initial weights; $\theta$, working copy of the next iterate}
\KwHyper{$B$, minibatch size; $\lambda$; $\beta_{\max}$; $M_{\tt safe}$}
\KwHyper{$\beta_{\tt fix}$, $k_{\tt warm}$, $M$, $\kappa_0$, $\kappa_{\min}$, $\kappa_{\max}$, $\gamma_\kappa$, $\eta_\kappa$; chosen number of gradient-descent steps per iterate}
\KwOut{$(\theta_k)_{k\ge0}$}
Set $\overline M_{-1}\gets M$\;
\For{$k=0,1,2,\dots$}{
  Define $\beta_k\gets\beta_{\tt fix}$ if $k<k_{\tt warm}$; otherwise define it from the frozen public heads $T^{\theta_k}$ and $C^{\theta_k}$, the current inverse multiplier $\kappa_k$, and the KKT rule above\;
  Sample $(\mathfrak S,\ell,N,\widetilde N,\mathfrak r)\gets
  \algcall{alg:sample-tree}{SampleTree}(\pi^{\theta_k},\beta_k;C^{\theta_k},M_{\tt safe})$\;
  Set $M_k^{\tt cf}\gets|\mathfrak S|+
  \sum_{\mathsf s\in\mathfrak S}(\widetilde N(\mathsf s)-N(\mathsf s))
  C^{\theta_k}(\ell(\mathsf s))$.  If $k\geq k_{\tt warm}$, set $\overline M_k\gets(1-\gamma_\kappa)\overline M_{k-1}+\gamma_\kappa M_k^{\tt cf}$ and update $\log\kappa_{k+1}$ by the controller above; otherwise set $(\overline M_k,\kappa_{k+1})\gets(\overline M_{k-1},\kappa_k)$\;
  Compute target table $\mathcal O\gets \algcall{alg:branched-value}{BranchedValue}(\mathfrak{S}, V^{\theta_k}, T^{\theta_k}, C^{\theta_k},\widetilde N,\mathfrak r)$\;
  \For{each gradient-descent step}{
    Sample $(P_b, \mathsf{S}_b)_{b=1}^B \stackrel{\mathsf{iid}}{\sim} (P, \mathsf{S})$ and set $\Sigma_b=\ell(\mathsf{S}_b)$\;
    Deduce $\widehat{A}_b^\lambda(x_{P_b}, \mathsf{a})$, $\widehat{Q}_b^\lambda(x_{P_b}, \mathsf{a})$, $\widehat{V}_b(x_{P_b})$, $\widehat T_b$ and $\widehat C_b^\lambda$ from $\mathcal O[\mathsf S_b].\mathrm{targets}$, and deduce $\widehat{\mathbb P}(x_{P_b}\mid\mathsf S_b)$ and $\widehat{\mathbb P}(x_{P_b}\mid\mathsf S_b\mathsf a)$ from the corresponding $\mathrm{posterior}$ fields\;
    Estimate
    \begin{equation}[rCl]
      \widehat{l}(\theta)
      &=& \frac{1}{B}\sum_{b=1}^B\sum_{\substack{x_{P_b} \\ \mathsf{a}\in\mathfrak{E}_{\mathsf{S}_b}}}
        \frac{1}{N(\mathsf{S}_b)}
        \mathcal{L}^\pi\left(
        (\pi^\theta, V^\theta,\pi^{\theta_k}, V^{\theta_k})(x_{P_b}, \Sigma_b),
        a_{\mathsf{a}},
        \widehat{A}_b^\lambda(x_{P_b}, \mathsf{a}),
        \widehat{Q}_b^\lambda(x_{P_b}, \mathsf{a})
        \right)
        \widehat{\mathbb{P}}(x_{P_b} \mid \mathsf{S}_b\mathsf a)
      \nonumber\\
      && + \frac{1}{B}\sum_{b=1}^B\sum_{x_{P_b}}
      \mathcal{L}^{V}\left(
        (V^\theta, V^{\theta_k})(x_{P_b}, \Sigma_b),
        \widehat{V}_b(x_{P_b})
        \right)\widehat{\mathbb{P}}(x_{P_b} \mid \mathsf{S}_b)
      \nonumber\\
      && + \frac{1}{B}\sum_{b=1}^B
      \mathcal L^{T,C}\left(
        (T^\theta,C^\theta)(\Sigma_b),
        \widehat T_b,
        \widehat C_b^\lambda
      \right)
        \label{eq:branched-minibatch-loss}
    \end{equation}\;
    Descend $\theta$ along $\nabla_\theta \widehat{l}(\theta)$\;
  }
  Store $\theta_{k+1} \gets \theta$\;
}
\Return{$(\theta_k)_{k\ge0}$}
\end{algorithm}

\section{Loss Estimator Variance}\label{sec:loss-estimator-variance}

\subsection{Sequential Rao--Blackwellization}\label{subsec:sequential-rb-loss-variance}

Freeze $\pi=\pi^{\theta_k}$, evaluate derivatives at $\theta=\theta_k$, and
take $\lambda=1$ with loss $\mathcal L_{\tt PG}+\mathcal L_{\tt V}$.  For a
player decision $\nu=(\sigma,a)$, put
\[
  \Phi_k(x,\tau;\nu)
  :=-\left(R_{o(\sigma)}(x,\tau)-V^{\theta_k}(x_{o(\sigma)},\sigma)\right)
  \left[\nabla_\theta\log\pi^\theta(a\mid x_{o(\sigma)},\sigma)
  +2c_V\nabla_\theta V^\theta(x_{o(\sigma)},\sigma)\right]_{\theta=\theta_k}.
\]
Thus conditioning $\Phi_k$ on the public pre-terminal trajectory gives the
exact Rao--Blackwellized gradient contribution.  Throughout this subsection,
$\operatorname{Var}(Y):=\operatorname{trCov}(Y)$.

Let $\mathfrak T^{\mathsf{pre-c}}=((\Tau_i^{\mathsf{pre-c}},\mathfrak r_i))_{i=1}^N$
be as in \cref{alg:sample-preterminal}, let
$D_i:=|\Tau_i^{\mathsf{pre-c}}|_{\textnormal{player}}$,
$L:=\sum_iD_i$, and $\ell:=L/N$.  Given the rollout batch, sample
$((I_b,P_b,\Sigma_b,A_b))_{b=1}^B$ from
\cref{def:preterminal-rao-blackwell-decision-law}.  For every rollout $i$, use
$n$ conditionally iid completions, with completion $j$ shared by all sampled
decisions for which $I_b=i$.  The scaled gradient estimator is
\[
  \widehat g_{N,B,n}
  :=\frac{\ell}{nB}\sum_{j=1}^n\sum_{b=1}^B
  \Phi_k(X_{I_b,j},\Tau_{I_b,j};\Sigma_b,A_b).
\]
The factor $\ell$ cancels the random denominator $L$ in the empirical decision
law, so $\widehat g_{N,B,n}$ is unbiased for the expected trajectory-sum
gradient.

For a generic public pre-terminal trajectory, write $D:=|\Tau^{\mathsf{pre-c}}|_{\textnormal{player}}$
and $\ell_\star:=\mathbb ED$.  For every player decision $\nu\preccurlyeq
\Tau^{\mathsf{pre-c}}$, define
\[
  \mu_\nu:=\mathbb E^\pi(\Phi_k(X,\Tau;\nu)\mid\Tau^{\mathsf{pre-c}}),
  \qquad
  G:=\sum_{\nu\preccurlyeq\Tau^{\mathsf{pre-c}}}\mu_\nu,
  \qquad
  Q:=\sum_{\nu\preccurlyeq\Tau^{\mathsf{pre-c}}}\|\mu_\nu\|_2^2,
\]
and $\xi_\nu:=\Phi_k(X,\Tau;\nu)-\mu_\nu$.  The constants governing rollout
and minibatch variance are
\begin{equation}[rCl]
  A_{\rm rollout}&:=&\operatorname{Var}(G),\nonumber\\
  A_{\rm MB}^{\ne}&:=&\ell_\star\mathbb EQ-\|\mathbb EG\|_2^2,\nonumber\\
  A_{\rm MB}^{=}&:=&\mathbb E(DQ-\|G\|_2^2),
  \label{eq:sequential-rb-rollout-minibatch-constants}
\end{equation}
while the completion constants are
\begin{equation}[rCl]
  C_{\rm MC}^{\ne}
  &:=&\ell_\star\mathbb E\sum_{\nu\preccurlyeq\Tau^{\mathsf{pre-c}}}
  \mathbb E^\pi(\|\xi_\nu\|_2^2\mid\Tau^{\mathsf{pre-c}}),\nonumber\\
  C_{\rm MC}^{=}
  &:=&\mathbb E\left[D\sum_{\nu\preccurlyeq\Tau^{\mathsf{pre-c}}}
  \mathbb E^\pi(\|\xi_\nu\|_2^2\mid\Tau^{\mathsf{pre-c}})\right],\nonumber\\
  C_{\rm MC}^{\rm pair}
  &:=&\mathbb E\left[\mathbb E^\pi\left(
  \left\|\sum_{\nu\preccurlyeq\Tau^{\mathsf{pre-c}}}\xi_\nu\right\|_2^2
  \middle|\Tau^{\mathsf{pre-c}}\right)\right].
  \label{eq:sequential-rb-montecarlo-constants}
\end{equation}
Here $A_{\rm rollout}$ is public-trajectory variance; the superscripts $\ne$
and $=$ distinguish different-rollout and same-rollout terms in the
second-moment expansion; and $C_{\rm MC}^{\rm pair}$ contains the covariance
between decisions sharing one completion.

Successive applications of the law of total variance give the exact
decomposition
\begin{equation}[rCl]
  \operatorname{Var}(\widehat g_{N,B,n})
  &=&T_{\rm rollout}(N)+T_{\rm minibatch}(N,B)+T_{\rm montecarlo}(N,B,n),\nonumber\\
  T_{\rm rollout}(N)
  &=&\dfrac{A_{\rm rollout}}{N},\nonumber\\
  T_{\rm minibatch}(N,B)
  &=&\dfrac{A_{\rm MB}^{\ne}}{B}
  +\dfrac{A_{\rm MB}^{=}-A_{\rm MB}^{\ne}}{BN},\nonumber\\
  T_{\rm montecarlo}(N,B,n)
  &=&\dfrac1n\left(
  \dfrac{C_{\rm MC}^{\ne}}{B}
  +\dfrac{C_{\rm MC}^{\rm pair}}{N}
  +\dfrac{C_{\rm MC}^{=}-C_{\rm MC}^{\ne}-C_{\rm MC}^{\rm pair}}{BN}
  \right).
  \label{eq:sequential-rb-loss-variance}
\end{equation}
The $C_{\rm MC}^{\rm pair}/(Nn)$ term is the variance floor caused by sharing
one completion across all selected decisions on a rollout; increasing $B$
alone cannot remove it.  All six constants are independent of $(N,B,n)$.

\section{Poker: a Perfect Example}

We focus on configurable equal-stack Texas Hold'em cash poker: $P\in\{2,\ldots,6\}$ players,
$S\in[8,128]$ big blinds per player in half-big-blind increments, and the highest
$R\in\{4,\ldots,13\}$ ranks with four suits. Let $D=4R$ and require $2P+5<D$.
There is no ante, straddle, or rake, and all legal actions in the fixed chip
denomination are allowed.
Below are a few core variables:
\begin{equation*}[rCl]
  |\mathcal{P}| &=& P\\
  \mathcal{P} &=& \{0,\ldots,P-1\}\\
  \mathcal{D}_R &=& \{52-D,\ldots, 51\}\\ % evaluator-facing suffix deck
  \mathcal{X}_p &=& \mathcal{H}_R = \binom{\mathcal{D}_R}{2}\\
  \mu &=& \mathsf{Unif}\left(\{x \in \mathcal{X} \mid x_p\cap x_q=\varnothing\textnormal{ for }p\ne q\}\right)\\
  \mathsf{B}(\sigma) &\in& \binom{\mathcal{D}_R}{0}\textnormal{ or }\binom{\mathcal{D}_R}{3}\textnormal{ or }\binom{\mathcal{D}_R}{4}\textnormal{ or }\binom{\mathcal{D}_R}{5}\\ %current board
  \mathsf{C}(x,\sigma) &=& \left(\bigcup_p x_p\right)\cup \mathsf{B}(\sigma)\\ %used cards
  \chi(\cdot \mid x, \sigma) &=& \mathsf{Unif}\left(\binom{\mathcal{D}_R \setminus \mathsf{C}(x,\sigma)}{k(\sigma)}\right),
  \qquad
  k(\sigma)=\begin{cases}[l"l]3&\textnormal{at the flop},\\1&\textnormal{at the turn or river}.\end{cases}
\end{equation*}

Several specificities impact our implementation:
\begin{enumerate}
  \item We feed the model \textbf{redundant} information by feeding it both the public state $\sigma$ and the Bayesian-updated ranges $\mathfrak{r}$.
  
  \item The game of poker has a strong \textbf{preflop gauge symmetry}: any suit permutation should leave the strategy invariant.
  Instead of $\binom{4R}{2}$ possible private player states, there are $R^2$
  classes (1326 and 169 respectively for the full deck).
  We exploit this by having a separate model for preflop and postflop, making it impossible for the model to not pick actions invariants under this transformation preflop.
  Optionally, we can also exploit this symmetry in the Monte Carlo estimators (see below) by storing the $4!$ suit permutations of each sampled private state and board into the accumulated statistics (this can be considered if possible with minimal wrap divergence).
  
  \item The \textbf{postflop gauge symmetry} is harder to handle as it acts on the board which is an input.
  Optionally, we can add a loss function term that: samples a random suit permutation, applies it to the model inputs, and penalizes the difference between the re-permutated outputs and the original outputs.
  This adds one forward pass per sample, which shouldn't be too expensive compared to the cost of a backward pass. 
  
  \item In \textbf{postflop} states, to control the size of the input, we actually only feed the ranges of the \textbf{live} players (meaning those who have not folded yet and are thus still eligible to win part of the pot).
  We feed the model information about folded players by exploiting the very architecture of attention heads (see subsection to come on cluster expansions and Möbius expansions).
  This adds a computational step which can be seen as the first step of the model (and can be launched in the same batches as the model itself).

  \item The model is a multi-headed attention transformer model, with one token per pair $(p,h)$ (for $p$ only a live player postlfop), one token for each piece of public information, and one taken per \textbf{action mode}.
  The tokens corresponding to action $a$ and pair $(p_{\textnormal{acting}, h})$ are used to output propensity for $a$ to be played with hand $h$.

  \item As a consequence, we need to either:
  \begin{enumerate}
    \item \textbf{Memory Pressure:} during rollout, store the live ranges of each rolled-out state, so they can be fed to the minibatch samples as inputs
    \item \textbf{Redundant Computations:} during minibatches, recompute the live ranges of each samples state $\Sigma$ by quickly walking back the tree from $\Sigma$ to $\varnothing$, and doing $\odot$-Bayesian updates along the way.
  \end{enumerate}

  \item Three operations are already implemented in \texttt{CUDA} and \texttt{C++} in a separate repository:
  \begin{enumerate}
    \item The \textbf{folded statistics} to be fused with the model as its first step.
    This is \textbf{only} used: right after a player folded (otherwise can be stored and passed down the trajectory), and postlfop (preflop we can afford feeding the model all ranges). 
    \item A \textbf{rejection sampler} to be used exclusively for the branching rollout, and to be fused with the model as its last step,
    \item A \textbf{Monte Carlo estimator} of equity and posteriors, used by both Rao--Blackwellized variants after rollouts are completed and before backpropagation to compute $\widehat{Q}$, $\widehat{V}$, \ldots
  \end{enumerate} 
  Folded statistics might be more efficiently computed by the CPU considering that they are only randomly necessary.
  The other two operations should really be GPU-based.

  \item We implement a synchronous rollout that is heavily computed by GPUs (typically a node of 2-4).
  We maintain a queue of states to be rolled out, and we roll them out in batches (two types of batches: one for preflop, one for postflop).

  \item \textbf{Research loss design:} the full design combines clipped \texttt{PPO} with an optional \texttt{KL} term to an older reference index $\operatorname{ref}(k)$, an entropy term, and losses for each proposed value head:
  \[
    \mathcal L_{\tt clip}[\widehat A^\lambda]+\mathcal L_{\tt H} + \mathcal L_{\tt KL} + \mathcal L_{\tt V}[\widehat V^\lambda] + \mathcal L^{T,C}[\widehat T,\widehat C^\lambda]
  \]
  The implementation companion exposes explicit \texttt{fixed} and
  \texttt{dynamic} branching modes. Both use the clipped policy, actor-type
  value, entropy, optional exact legal-action reference-policy KL, and public
  $T/C$ losses from the first update. The dynamic mode follows
  $\beta_{\tt fix}$ during warmup and only afterward lets the KKT controller
  consume $T/C$ for rollout control. \texttt{NeuRD} and stored all-action $Q$
  targets remain deferred research extensions.

  \item Ideally, the implemtation should be clearly modular: it will be used for poker, but must be easily adaptable to other PBGs.
\end{enumerate}

\end{document}
